% !TeX root = ../../Book2.tex
% 纲要标题：## 第10章　双线性型、二次型与经典群
% 纲要位置：计划纲要.md，第 1372 行。

\chapter{双线性型、二次型与经典群}
\label{b2:ch:10}
\BookChapterNumber{b2:ch:10}

\begin{chapterabstract}
一个型不仅给向量赋值，还规定哪些方向相互正交。双线性型的换基产生合同，而线性算子的换基产生相似；本章先分清这两种分类问题，再建立交替型标准形、特征不为二的二次型对角化及实惯性定律。一般域上的二次型则通过双曲平面的分离、反射与 Witt 消去整理为双曲部分和各向异性部分。最后讨论保持各种型的经典群，并以 Clifford 代数的商构造说明二次型与结合代数的联系。
\end{chapterabstract}

\begin{learninggoals}
\begin{itemize}
\item 辨认非退化配对、根空间与正交补，说明何时正交补给出直和分解。
\item 对具体矩阵作合同化简，比较不同域及不同特征下的二次型。
\item 掌握双曲平面的分离和反射，能够逐步说明 Witt 分解的存在性、唯一性与消去。
\item 从保持型的条件写出正交群、辛群与酉群，区分双线性与 Hermitian 约定，并理解 Clifford 代数的泛性质。
\end{itemize}
\end{learninggoals}

在实平面上，\(q_+(x,y)=x^2+y^2\) 与 \(q_-(x,y)=x^2-y^2\) 的矩阵都可逆，但两者的几何行为不同：前者只在原点为零，后者在直线 \(x=y\) 和 \(x=-y\) 上也为零。若把底域改为复数，坐标代换 \(y'=iy\) 就有 \(x^2-y^2=x^2+(y')^2\)，两型因而等距。一个型的分类因此不仅要问矩阵是否可逆，还要问在什么域上、允许怎样换基，以及哪些向量与子空间被它判为正交或零值。

\ProseParagraphBreak

对称型可以通过对角化逐步拆解；交替型和一般双线性型却需要不同的操作。特别在特征二，\(q(v+w)-q(v)-q(w)\) 仍给出极化信息，却不能靠除以 \(2\) 从中恢复 \(q\)。把这些适用条件说清，再看保持型的变换，正交群及由二次关系构造的代数才有准确的起点。

\ProseParagraphBreak

本章需要\chapref{b2:ch:01}的有限维向量空间、对偶、矩阵换基及行列式的基本性质。\secref{b2:sec:1001}对任意域讨论双线性型；\secref{b2:sec:1002}明确区分特征二；\secref{b2:sec:1003}始终假设特征不为二，并在有限维非退化空间中证明 Witt 理论。\subsecref{b2:subsec:1004:03}还使用\cref{b2:thm:tensor-algebra-universal}的张量代数泛性质。

\ProseParagraphBreak
双线性型及其几何可参阅 Roman 的《Advanced Linear Algebra》\cite[Chapter 11]{Roman2008AdvancedLinearAlgebra}；实惯性与 Witt 消去分别见同书\cite[Theorems 11.25--11.26, 11.34]{Roman2008AdvancedLinearAlgebra}。比较文献时应先检查极化型是否除以二。本章在特征不为二时统一写 \(q(v)=b(v,v)\)，因而双曲平面的表达式为 \(2xy\)。

% !TeX root = ../../Book2.tex
% 纲要标题：### 10.1 双线性型
% 纲要位置：计划纲要.md，第 1374 行。

\section{双线性型}
\label{b2:sec:1001}

线性映射接受一个向量，双线性型则同时比较两个向量。选定基以后，两者都能写成矩阵，但换基时的规则不同。本节先把配对两端的空间分开；只有在说明对称性或交替性以后，才把左右正交条件合并。

% 纲要要点（供目录核对）：
%
% * 配对、有限维非退化性及正交补；用对偶限制映射判定维数与直和
% * 对称型与交替型；交替型的二维分离在任意特征成立
% * 型的合同与算子的相似
%

\subsection{配对、非退化性与正交补}
\label{b2:subsec:1001:01}

\begin{definition}\label{b2:def:bilinear-pairing}
设 \(k\) 为域，\(U,V\) 为 \(k\) 向量空间。如果映射 \(b:U\times V\to k\) 在固定任意一个变量后对另一个变量都是 \(k\) 线性的，就称 \(b\) 为\term[shuangxianxing-peidui]{双线性配对}[bilinear pairing]。它的左根空间与右根空间分别为
\[
\{\,u\in U\mid b(u,v)=0\;\text{对所有}\;v\in V\,\},\qquad
\{\,v\in V\mid b(u,v)=0\;\text{对所有}\;u\in U\,\}.
\]
两者均为零时，称配对\term[peidui-feituihua]{非退化}[nondegenerate]。当 \(U=V\) 时，称 \(b\) 为 \(V\) 上的\term[shuangxianxing-xing]{双线性型}[bilinear form]。
\end{definition}

以下空间均有限维。固定 \(u\in U\) 后，它与第二侧所有向量的配对值组成 \(V\) 上的线性泛函 \(b(u,-)\)；固定 \(v\in V\) 则得到 \(U\) 上的线性泛函 \(b(-,v)\)。因此配对把一侧的向量变成另一侧的线性泛函，给出两个映射
\[
L_b:U\longrightarrow V^*,\quad u\longmapsto b(u,-),
\qquad
R_b:V\longrightarrow U^*,\quad v\longmapsto b(-,v).
\]
它们的核就是上述两个根空间。对这样的有限维 \(k\) 向量空间上的配对 \(b:U\times V\to k\)，若 \(L_b:U\to V^*\) 是同构，就称 \(b\) 为\term[wanmei-peidui]{完美配对}[perfect pairing]；这表示 \(V\) 上的每个线性泛函都唯一地由某个 \(u\in U\) 与 \(V\) 配对得到。

\begin{proposition}\label{b2:prop:finite-perfect-pairing}
设 \(k\) 为域，\(U,V\) 为有限维 \(k\) 向量空间，\(b:U\times V\to k\) 为双线性配对。令 \(L_b:U\to V^*\)、\(u\mapsto b(u,-)\) 及 \(R_b:V\to U^*\)、\(v\mapsto b(-,v)\)，则 \(b\) 非退化当且仅当 \(L_b\) 与 \(R_b\) 都是同构；这也等价于其中任一个是同构。此时 \(\dim U=\dim V\)。
\end{proposition}
\begin{proof}
非退化使两个映射都单射，故 \(\dim U\le\dim V\le\dim U\)，它们因而同为同构。若只知 \(L_b\) 是同构，且 \(R_b(v)=0\)，则全部 \(U\) 中元素都与 \(v\) 配为零。\(L_b\) 满射说明每个 \(V\) 上的线性泛函都在 \(v\) 处为零，由\cref{b2:lem:functional-extension}的分离结论，\(v=0\)。又两空间等维，所以 \(R_b\) 也是同构。交换两个变量证明另一方向。
\end{proof}

例如，求值 \(V^*\times V\to k\) 在有限维时完美。若取 \(b:k^2\times k^3\to k\)、\(b(u,v)=u_1v_1+u_2v_2\)，左根空间为零，右根空间却是一维；只检验一端不够。

\ProseParagraphBreak
上面从两个单射推出两个同构，使用了有限维的维数比较。若取无限维空间 \(V=k^{(\mathbb N)}\) 及配对 \(b(u,v)=\sum_i u_iv_i\)，每个向量支集有限使求和有限，且用标准基向量 \(e_i\) 检测便知两端根空间都为零。可是线性泛函 \(\ell(v)=\sum_i v_i\) 满足 \(\ell(e_i)=1\) 对每个 \(i\) 都成立，不可能等于任何 \(b(u,-)\)，因为 \(u\) 的支集有限；因此这个非退化配对的 \(L_b\) 并不满射。

\ProseParagraphBreak
对于同一空间上的一般型，条件 \(b(u,v)=0\) 与 \(b(v,u)=0\) 未必等价。以下两种型允许我们合并这些条件。

\begin{definition}\label{b2:def:symmetric-alternating-form}
设 \(k\) 为域，\(V\) 为 \(k\) 向量空间，\(b:V\times V\to k\) 为双线性型。如果 \(b(u,v)=b(v,u)\) 对所有 \(u,v\in V\) 成立，就称 \(b\) 为\term[duichen-xing]{对称型}[symmetric bilinear form]；如果 \(b(v,v)=0\) 对所有 \(v\in V\) 成立，就称它为\term[jiaoti-xing]{交替型}[alternating bilinear form]。如果 \(b(u,v)=-b(v,u)\) 对所有 \(u,v\in V\) 成立，就称它为\term[fanduichen-xing]{反对称型}[skew-symmetric bilinear form]。
\end{definition}

将 \(b(u+v,u+v)=0\) 展开，便知交替型总是反对称的。对于子空间 \(W\)，先分别记
\[
W^{\perp_r}=\{\,v\in V\mid b(w,v)=0\ (w\in W)\,\},\qquad
{}^{\perp_l}W=\{\,v\in V\mid b(v,w)=0\ (w\in W)\,\}.
\]
若 \(b\) 对称，或满足 \(b(v,u)=-b(u,v)\)，两者相同，统记为 \(W^\perp\)，称为\term[zhengjiao-bu]{正交补}[orthogonal complement]；此时 \(\operatorname{rad}b=V^\perp\) 称为型的\term[xing-de-gen]{根空间}[radical of a form]。这里的“补”首先是正交条件的名称，并不预先断言 \(V=W\oplus W^\perp\)。
\symidx{\(W^\perp\)：对称型或交替型下的正交补}
\symidx{\(\operatorname{rad}b\)：双线性型的根空间}

\begin{theorem}\label{b2:thm:orthogonal-complement}
设 \(k\) 为域，\(V\) 为有限维 \(k\) 向量空间，\(b:V\times V\to k\) 为对称型或交替型，\(W\le V\)。记 \(W^\perp=\{v\in V\mid b(v,w)=0\;\text{对所有}\;w\in W\}\)。若 \(b\) 在 \(V\) 上非退化，则
\[
\dim W^\perp=\dim V-\dim W,\qquad (W^\perp)^\perp=W.
\]
无论 \(b\) 在 \(V\) 上是否退化，下述正交直和分解成立当且仅当 \(b|_{W\times W}\) 非退化：
\[
V=W\mathbin{\perp}W^\perp,
\]
其中符号 \(\perp\) 表示两个分量互相正交的直和。在上述分解成立的前提下，若 \(b\) 在 \(V\) 上也非退化，则 \(W^\perp\) 上的限制型非退化；若 \(b\) 在 \(V\) 上退化，则不保证这一点。
\end{theorem}
\begin{proof}
正交补是满足全部配对方程的向量组成的核。具体地，考虑线性映射
\[
\Phi:V\longrightarrow W^*,\qquad
\Phi(v)(w)=b(w,v)\quad(w\in W).
\]
由型的对称性或反对称性，\(\ker\Phi=W^\perp\)。当 \(V\) 上非退化时，\(\Phi\) 满射：给定 \(W\) 上的泛函，先由\cref{b2:lem:functional-extension}延拓到 \(V\)，再由\cref{b2:prop:finite-perfect-pairing}写成 \(b(-,v)\)。于是秩与核的维数公式给出第一个等式。正交条件的对称性给出 \(W\subseteq(W^\perp)^\perp\)，再由维数相等得到相等。

限制型的根空间满足
\[
\operatorname{rad}(b|_{W\times W})=W\cap W^\perp,
\]
因为两边都由 \(W\) 中与全部 \(W\) 配为零的向量组成。因此正交直和分解若成立，限制型就非退化。反过来，现在只假设 \(W\) 上非退化。对任意 \(v\in V\)，\cref{b2:prop:finite-perfect-pairing}在 \(W\) 上给出唯一 \(w\in W\)，使 \(b(x,w)=b(x,v)\) 对所有 \(x\in W\) 成立。因此 \(v-w\in W^\perp\)，而上述根空间等式给出 \(W\cap W^\perp=0\)，这就证明直和分解。若 \(V\) 上也非退化，而 \(z\in W^\perp\) 与全部 \(W^\perp\) 正交，则它还与 \(W\) 正交，故与 \(V\) 正交，只能为零。
\end{proof}

若 \(b\) 在 \(k^2\) 上的矩阵为 \(\operatorname{diag}(1,0)\)，取 \(W=ke_1\)，则 \(W\) 上非退化，且 \(W^\perp=ke_2\)；后者上的限制型却为零。

\ProseParagraphBreak
例如，实平面上 \(b(x,y)=x_1y_1-x_2y_2\) 非退化，但直线 \(W=\mathbb R(1,1)\) 满足 \(W^\perp=W\)。维数公式正确，直和分解却不成立，因为限制型为零。非退化性不能从整个空间自动传给每个子空间。

\subsection{对称型与交替型}
\label{b2:subsec:1001:02}

交替性已经蕴含反对称性。特征不为二时，反对称性反过来给出 \(2b(v,v)=0\)，所以二者等价。特征二时不能这样约去二；对称型 \(b(x,y)=xy\) 就不是交替型。交替型的定义在任意特征下都要求对角取值为零。

\begin{theorem}[交替型的标准形]\label{b2:thm:alternating-form-normal}
设 \(k\) 为任意特征的域，\(V\) 为有限维 \(k\) 向量空间，\(b:V\times V\to k\) 为交替双线性型，则 \(V\) 有一组基
\[
e_1,f_1,\ldots,e_m,f_m,z_1,\ldots,z_r,
\]
使不同的两维分量及最后的根空间互相正交，并有
\[
b(e_i,f_i)=1,\quad b(f_i,e_i)=-1,\quad b(e_i,e_i)=b(f_i,f_i)=0.
\]
\(z_1,\ldots,z_r\) 是根空间的基。因此型的秩为 \(2m\)，并完全决定其在给定维数下的合同类型。特别地，非退化交替型只存在于偶数维空间。
\end{theorem}
\begin{proof}
若型为零，任取基即可。否则要分离一个非退化的两维平面，再在其正交补中继续同样的构造；交替型的每条直线上的限制型都为零，所以这里不能只分离一条直线。选 \(e,f_0\) 使 \(b(e,f_0)\ne0\)，将 \(f_0\) 除以该值，得到 \(b(e,f)=1\)。\(e,f\) 线性无关：若 \(ae+cf=0\)，分别与 \(e,f\) 配对便得 \(c=a=0\)。它们的线性生成空间 \(H\) 上的矩阵为
\[
\begin{pmatrix}0&1\\-1&0\end{pmatrix},
\]
行列式为一，故限制型非退化。由\cref{b2:thm:orthogonal-complement}中不要求整体非退化的分解，\(V=H\perp H^\perp\)。在较低维的 \(H^\perp\) 上归纳，直到剩余限制型为零。每个已分离的两维块非退化，故最后的零型分量恰是整个型的根空间。

上述基给出的矩阵由 \(m\) 个同样的可逆两维块及一个 \(r\) 维零块组成，所以秩为 \(2m\)。秩不随换基改变，\(r=\dim V-2m\) 也被确定；同秩的两个型把上述基逐一对应便得到保持配对的同构。
\end{proof}

非退化交替型的一组 \(e_i,f_i\) 称为\term[xin-ji]{辛基}[symplectic basis]。若先排列 \(e_1,\ldots,e_m\)，再排列 \(f_1,\ldots,f_m\)，其矩阵就是
\[
J_m=\begin{pmatrix}0&I_m\\-I_m&0\end{pmatrix}.
\]
这些论证没有除以二，因而在特征二中仍然有效；此时 \(-I_m=I_m\)，单个辛块变成
\[
\begin{pmatrix}0&1\\1&0\end{pmatrix}.
\]
它既对称又交替：对称性来自 \(-1=1\)，交替性仍由对角线为零保证。

\subsection{型的合同与算子的相似}
\label{b2:subsec:1001:03}

选基 \(\mathcal E=(e_1,\ldots,e_n)\)，令 \(B=\bigl(b(e_i,e_j)\bigr)\)。若 \(x,y\) 是两个向量的坐标列，则
\[
b(u,v)=x^{\mathsf T}By.
\]
这既说明矩阵由型唯一确定，也说明每个矩阵都确定一个双线性型。非退化恰好等于 \(B\) 可逆；对称性等于 \(B^{\mathsf T}=B\)，交替性等于 \(B^{\mathsf T}=-B\) 且对角元全零。

\begin{proposition}\label{b2:prop:bilinear-congruence}
设 \(k\) 为域，\(V\) 为 \(n\) 维 \(k\) 向量空间，\(b:V\times V\to k\) 为双线性型，其旧基矩阵为 \(B\)。若新基的列在旧基中的坐标组成 \(P\in\operatorname{GL}_n(k)\)，则型的新矩阵为
\[
B'=P^{\mathsf T}BP.
\]
存在这样的 \(P\) 时，称两矩阵\term[hetong]{合同}[congruent]。两个型之间的\term[dengju-tonggou]{等距同构}[isometry]是保持双线性取值的线性双射；它们等距，当且仅当相应矩阵合同。
\end{proposition}
\begin{proof}
新坐标为 \(x',y'\) 时，旧坐标为 \(Px',Py'\)，故取值是 \((x')^{\mathsf T}P^{\mathsf T}BPy'\)。逐个代入标准基向量便得到矩阵公式。一般线性双射也由可逆矩阵给出，保持两变量取值的等式正是同一个合同等式。
\end{proof}

算子换基的公式则为 \(A'=P^{-1}AP\)：输入的新坐标先乘 \(P\) 变为旧坐标，作用 \(A\) 后，输出的旧坐标再乘 \(P^{-1}\) 变回新坐标。双线性型的两个位置都是输入变量，两侧的坐标都乘 \(P\)，所以给出 \(P^{\mathsf T}BP\)。实一维型的矩阵 \((1)\) 与 \((4)\) 通过 \(P=(2)\) 合同，但一维矩阵相似只能使其不变，故二者不相似。

\ProseParagraphBreak
合同保持秩。对于非退化型，等式
\[
\det B'=(\det P)^2\det B
\]
还说明行列式在 \(k^\times/(k^\times)^2\) 中的平方类不变。不过秩与行列式平方类合在一起仍不足以分类。例如实数域上 \(I_2\) 与 \(-I_2\) 都有秩二和行列式一，但在非零向量上，一个二次取值恒正，另一个恒负，所以不可能合同。下一节将从二次取值出发说明这里还缺少什么信息。

% !TeX root = ../../Book2.tex
% 纲要标题：### 10.2 二次型
% 纲要位置：计划纲要.md，第 1380 行。

\section{二次型}
\label{b2:sec:1002}

双线性型在两个位置取同一个向量，便得到二次表达式。这个过程能否逆转，取决于系数域的特征。本节始终令 \(V\) 有限维，并先在任意域上定义二次型；讨论对角化和惯性时，再明确加入特征与次序条件。

% 纲要要点（供目录核对）：
%
% * 特征不为2时的对角化及一维型的平方类判据
% * 实数域上的惯性定律；正负指数由定号子空间的最大维数确定
% * 特征2时二次型与极化型的区别；相同非退化极化型不足以判定等距
%

\subsection{特征不为二时的对角化}
\label{b2:subsec:1002:01}

\begin{definition}\label{b2:def:quadratic-form}
设 \(k\) 为域，\(V\) 为 \(k\) 向量空间。如果映射 \(q:V\to k\) 满足 \(q(av)=a^2\cdot q(v)\) 对所有 \(a\in k,v\in V\) 成立，且映射
\begin{equation}\label{b2:eq:quadratic-polar}
\beta_q(u,v)=q(u+v)-q(u)-q(v)
\end{equation}
是 \(k\) 双线性的，就称 \(q\) 为 \(V\) 上的\term[erci-xing]{二次型}[quadratic form]，称 \(\beta_q\) 为它的\term[jihua-xing]{极化型}[polar form]。带有指定二次型 \(q\) 的向量空间 \((V,q)\) 称为\term[erci-kongjian]{二次空间}[quadratic space]；这一名称允许 \(q\) 退化。
\end{definition}

第一个条件说明沿同一方向伸缩时，取值按标量的平方变化；极化型的双线性则控制不同方向之间的混合项。若 \(e_1,\ldots,e_n\) 是一组基，反复展开 \(q(u+v)=q(u)+q(v)+\beta_q(u,v)\)，再用缩放条件，就得到任意特征下都成立的坐标公式
\[
q\left(\sum_i x_ie_i\right)
=\sum_i x_i^2\cdot q(e_i)
+\sum_{i<j}x_ix_j\cdot\beta_q(e_i,e_j).
\]
因此这两个公理恰好使坐标表达式由平方项与两个不同坐标的乘积组成。

\ProseParagraphBreak
极化型总对称。当 \(\operatorname{char}k\ne2\) 时，本章采用归一化记号
\begin{equation}\label{b2:eq:normalized-polar}
b_q(u,v)=\frac12\beta_q(u,v),\qquad q(v)=b_q(v,v).
\end{equation}
第二个等式来自 \(\beta_q(v,v)=q(2v)-2q(v)=2q(v)\)：除以二正是为了使两个变量取同一向量时恢复 \(q\) 的值。反过来，对任意对称双线性型 \(b\)，令 \(q(v)=b(v,v)\)，则 \(\beta_q=2b\)。所以在特征不为二时，两类对象一一对应。以下在这一特征假设下，未加下标的 \(b\) 表示 \(b_q\)，未归一化的极化型仍记为 \(\beta_q\)。
\symidx{\(\beta_q\)、\(b_q=\beta_q/2\)：极化型及特征不为二时的归一化型}

在这项特征假设下，称 \(q\) 非退化，是指 \(b_q\) 非退化；两个二次型等距，是指线性双射保持其二次取值。由极化公式，这等价于保持 \(b_q\)。

\begin{theorem}[二次型的对角化]\label{b2:thm:quadratic-diagonalization}
设 \(k\) 为特征不为二的域，\(V\) 为有限维 \(k\) 向量空间，\(q:V\to k\) 为二次型，并令 \(b_q=\beta_q/2\)，则 \(V\) 有一组基，使
\[
q(x_1e_1+\cdots+x_ne_n)=a_1x_1^2+\cdots+a_rx_r^2,
\qquad a_i\ne0.
\]
最后 \(n-r\) 个基向量组成 \(\operatorname{rad}b_q\) 的基。特别地，每个对称矩阵都合同于对角矩阵。记 \(\langle a\rangle\) 为一维二次型 \(t\mapsto at^2\)。对任意 \(a,b\in k^\times\)，在上述 \(\operatorname{char}k\ne2\) 假设下，有
\[
\langle a\rangle\cong\langle b\rangle
\quad\Longleftrightarrow\quad
a/b\in(k^\times)^2,
\]
其中同构指保持二次取值的线性同构。
\end{theorem}
\begin{proof}
若 \(q\) 为零，极化公式使 \(b_q=0\)，任意基都合适。否则取 \(e\) 使 \(a=q(e)\ne0\)。直线 \(ke\) 上的限制型非退化，\cref{b2:thm:orthogonal-complement}给出
\[
V=ke\perp(ke)^\perp.
\]
为消去 \(e\) 方向与其余方向之间的混合项，要从 \(v\) 中减去一个 \(e\) 的倍数，使剩余向量与 \(e\) 正交。由 \(b(e,e)=a\)，这一倍数只能是 \(b(e,v)/a\)，因而任意 \(v\) 分解为
\[
v=\frac{b(e,v)}a e+\left(v-\frac{b(e,v)}a e\right),
\]
第二项与 \(e\) 正交。在较小维的正交补上归纳，得到全部非零一维分量及最后的零型分量。不同分量之间的混合项为零，所以二次表达式具有所述形式；从 \(b_q\) 的对角矩阵可见，最后这些基向量恰好生成 \(\operatorname{rad}b_q\)。

一维线性双射必为 \(t\mapsto ct\)，其中 \(c\in k^\times\)。
从 \(\langle a\rangle\) 到 \(\langle b\rangle\) 的这个映射保持二次取值，当且仅当 \(a=bc^2\)。
故两型等距当且仅当 \(a/b\) 是非零平方。
\end{proof}

常将非退化对角型写为 \(\langle a_1,\ldots,a_r\rangle\)。对角化并不表示这列系数已经唯一：由一维型的平方类判据，\(\langle1\rangle\) 与 \(\langle4\rangle\) 在 \(\mathbb Q\) 上等距，但与 \(\langle2\rangle\) 不等距。
\symidx{\(\langle a_1,\ldots,a_r\rangle\)：对角二次型}

\ProseParagraphBreak
例如，在 \(\mathbb Q\) 上取 \(q(x,y)=x^2+2xy+3y^2\)。令 \(u=x+y,v=y\)，便得 \(q=u^2+2v^2\)。新基在旧坐标中是 \((1,0),(-1,1)\)，两者对于矩阵 \(\begin{psmallmatrix}1&1\\1&3\end{psmallmatrix}\) 正交；新矩阵为 \(\operatorname{diag}(1,2)\)。这是同时变换两个变量的合同计算。

\begin{corollary}\label{b2:cor:quadratic-algebraically-closed}
设 \(k\) 为特征不为二的代数闭域，\(V\) 为有限维 \(k\) 向量空间，则 \(V\) 上的二次型在等距同构意义下由其秩完全分类。
\end{corollary}
\begin{proof}
由\cref{b2:thm:quadratic-diagonalization}对角化。每个非零 \(a_i\) 都有平方根，缩放相应基向量后将它改为一，得到 \(\operatorname{diag}(I_r,0)\)。反过来合同保持秩，故不同的 \(r\) 不等距。
\end{proof}

在代数闭域上，所有非零系数属于同一个平方类，所以对角化之后只剩秩；在 \(\mathbb Q\) 上，非零平方类已经有很多种，例如一与二不能彼此缩放得到。实数域上的非零系数则可缩为正一或负一，但这两个平方类不能合并。

\subsection{实数域上的惯性定律}
\label{b2:subsec:1002:02}

在 \(k=\mathbb R\) 上，非零平方不能改变符号。对角化后可以把每个非零系数的绝对值缩放为一，正负号却必须保留。对于实向量空间 \(V\) 上的二次型 \(q\)，若每个非零 \(v\in V\) 都满足 \(q(v)>0\)，就称 \(q\)\term[zhengding]{正定}[positive definite]；若 \(-q\) 正定，就称 \(q\) 负定。

\begin{theorem}[Sylvester 惯性定律]\label{b2:thm:sylvester-inertia}
设 \(V\) 为有限维实向量空间，\(q:V\to\mathbb R\) 为二次型，则 \(V\) 有一组基，使 \(q\) 在相应坐标下有标准表达式
\[
q(x)=x_1^2+\cdots+x_{p_+}^2-x_{p_++1}^2-\cdots-x_{p_++p_-}^2,
\qquad p_++p_-+r=\dim V.
\]
三元组 \((p_+,p_-,r)\) 唯一，且
\[
\begin{aligned}
p_+&=\max\{\,\dim W\mid W\le V,\ q|_W\text{ 正定}\,\},\\
p_-&=\max\{\,\dim W\mid W\le V,\ -q|_W\text{ 正定}\,\}.
\end{aligned}
\]
两个实二次型等距，当且仅当这三个数相同。
\end{theorem}
\begin{proof}
存在性由\cref{b2:thm:quadratic-diagonalization}及基向量乘 \(1/\sqrt{|a_i|}\) 得到。为证唯一性，记正项、负项及零项的坐标子空间为 \(V_+,V_-,V_0\)。\(V_+\) 上的型正定，维数为 \(p_+\)。若 \(W\) 是任意一个正定子空间，投影 \(W\to V_+\) 必单射：投影为零的 \(w\) 属于 \(V_-\oplus V_0\)，因而 \(q(w)\le0\)，正定性使 \(w=0\)。所以 \(\dim W\le p_+\)。这说明 \(p_+\) 是正定子空间的最大维数，与所选标准基无关。

对 \(-q\) 作同样论证，负项数 \(p_-\) 也由型确定。\(r=\dim V-p_+-p_-\) 因而确定；它也等于根空间的维数。相同三元组的型把各类标准基逐一对应，便得到等距同构，证明分类断言。
\end{proof}

这一结果称为\term[Sylvester-guanxing-dinglv]{Sylvester 惯性定律}[Sylvester's law of inertia]。\footnote{西尔维斯特（James Joseph Sylvester，1814-1897），英国数学家，研究矩阵、行列式与不变量理论。}正项数、负项数分别称\term[zheng-guanxing-zhishu]{正惯性指数}[positive index of inertia]、\term[fu-guanxing-zhishu]{负惯性指数}[negative index of inertia]，差 \(p_+-p_-\) 称为\term[fuhaocha]{符号差}[signature]。正定比非退化强。\(q(v)=0\) 只要求一个二次取值为零，而 \(v\in\operatorname{rad}b_q\) 要求 \(b_q(v,w)=0\) 对所有 \(w\) 成立；在特征不为二时，后者取 \(w=v\) 就推出前者，反向却不成立。例如实型 \(q(x,y)=x^2-y^2\) 非退化，\(v=(1,1)\) 满足 \(q(v)=0\)，但 \(b_q(v,(1,0))=1\)，所以不在根空间中。

\ProseParagraphBreak
例如 \(q(x,y,z)=2xy+z^2\)。令 \(u=(x+y)/\sqrt2,v=(x-y)/\sqrt2\)，得到 \(q=u^2-v^2+z^2\)，惯性为 \((2,1,0)\)。

\subsection{特征二的二次型与极化型}
\label{b2:subsec:1002:03}

现在假设 \(\operatorname{char}k=2\)。仍用\eqref{b2:eq:quadratic-polar}定义 \(\beta_q\)，但不能定义 \(\beta_q/2\)。又因 \(q(0)=0\)，有 \(\beta_q(v,v)=0\)，所以极化型必为交替型。

\begin{proposition}\label{b2:prop:quadratic-char-two-coordinates}
设 \(k\) 为特征二的域，\(V\) 为以 \(e_1,\ldots,e_n\) 为基的有限维 \(k\) 向量空间，\(q:V\to k\) 为二次型，\(\beta_q\) 为其极化型。在这组基的二次表达式中，混合项 \(x_ix_j\) 的系数为 \(\beta_q(e_i,e_j)\)（\(i<j\)）；各平方项 \(x_i^2\) 的系数 \(q(e_i)\) 可以任意改变而不改变极化型。因此，当 \(n\ge1\) 时，单凭 \(\beta_q\) 不能恢复 \(q\)。
\end{proposition}
\begin{proof}
坐标公式已由二次型的两个公理展开得到，混合项系数正是 \(\beta_q(e_i,e_j)\)。在特征二中，每个平方项的极化为零，因为 \((x+y)^2=x^2+y^2\)；改变它的系数仍得到二次型，且不改变极化型。若 \(n\ge1\)，改变一个平方项系数就改变相应 \(e_i\) 处的取值，所以极化型不能确定 \(q\)。
\end{proof}

例如一维上的零型与 \(q(x)=x^2\) 具有相同的零极化型，但不是同一个二次型。在 \(\mathbb F_2^2\) 上，\(q(x,y)=xy\) 的极化型为
\[
\beta_q\bigl((x,y),(u,v)\bigr)=xv+yu,
\]
其矩阵 \(\begin{psmallmatrix}0&1\\1&0\end{psmallmatrix}\) 可逆。特征二中，线性组合的平方仍是各平方项的线性组合，例如 \((ax+by)^2=a^2x^2+b^2y^2\)，所以换基也不能让平方项产生所需的混合项。任何仅含平方项的表达式 \(\sum a_ix_i^2\) 在任意基下都有零极化型，与这个 \(q\) 的非退化极化型矛盾。因此对称双线性型的矩阵分类不能替代特征二二次型的分类。

\ProseParagraphBreak
特征二时，映射 \(q\mapsto\beta_q\) 已经丢失平方项的信息；而特征不为二时，\(q(v)=\dfrac12\beta_q(v,v)\) 能恢复全部取值。“极化型非退化”“二次型有非零零点”与“根空间非零”需要分别说明，尤其不能由极化型为零就断言二次型为零。即使极化型非退化，也可能有两个不等距的二次型与它对应。

\begin{proposition}\label{b2:prop:quadratic-char-two-polar-failure}
设 \(V=\mathbb F_2^2\)，定义二次型
\[
q_0(x,y)=xy,\qquad q_1(x,y)=x^2+xy+y^2.
\]
两者有相同的非退化极化型
\[
\beta_{q_0}\bigl((x,y),(u,v)\bigr)
=\beta_{q_1}\bigl((x,y),(u,v)\bigr)=xv+yu,
\]
但二次空间 \((V,q_0)\) 与 \((V,q_1)\) 不等距。
\end{proposition}
\begin{proof}
两个表达式均按标量的平方缩放；直接展开极化公式，平方项在特征二中消去，两者都得到 \(xv+yu\)，所以它们确为二次型。这个极化型的矩阵为 \(\begin{psmallmatrix}0&1\\1&0\end{psmallmatrix}\)，可逆，因而非退化。

空间恰有三个非零向量 \((1,0),(0,1),(1,1)\)。逐个计算得
\[
\begin{aligned}
q_0(1,0)&=0,&q_0(0,1)&=0,&q_0(1,1)&=1,\\
q_1(1,0)&=1,&q_1(0,1)&=1,&q_1(1,1)&=1.
\end{aligned}
\]
因此 \(q_0\) 有两个非零零点，\(q_1\) 没有非零零点。线性双射会置换这三个非零向量，而等距同构还必须保持二次取值，所以它不可能把 \(q_0\) 变成 \(q_1\)。
\end{proof}

下一节讨论的 Witt 分解重新假设特征不为二。

% !TeX root = ../../Book2.tex
% 纲要标题：### 10.3 Witt 理论初步
% 纲要位置：计划纲要.md，第 1386 行。

\section{Witt 理论初步}
\label{b2:sec:1003}

从本节起固定 \(\operatorname{char}k\ne2\)，所讨论的二次空间均有限维，且以 \(b=\beta_q/2\) 为相应对称型。实数域上的正负号分类不能直接搬到一般域上，但非零零点仍能帮助我们分离出一种固定的两维型。Witt 理论先取出这些两维部分，再研究剩下没有非零零点的部分。

% 纲要要点（供目录核对）：
%
% * 双曲平面、各向同性与各向异性；各向同性不等于属于根空间
% * Witt 分解与消去定理；反射用于对齐共同非退化分量
% * Witt 群的概念；忽略双曲部分后由相反型给出加法逆元
%

\subsection{双曲平面、各向同性与各向异性}
\label{b2:subsec:1003:01}

\begin{definition}\label{b2:def:isotropic-hyperbolic}
设 \(k\) 为特征不为二的域，\((V,q)\) 为 \(k\) 上的二次空间。若 \(v\in V\) 满足 \(v\ne0\) 且 \(q(v)=0\)，就称 \(v\) 为\term[gexiang-tongxing-xiangliang]{各向同性向量}[isotropic vector]。若不存在这样的向量，就称 \((V,q)\)\term[gexiang-yixing]{各向异性}[anisotropic]；零空间按此定义也是各向异性的。若子空间 \(T\subseteq V\) 满足 \(q|_T=0\)，就称 \(T\) 为\term[quan-gexiang-tongxing]{全各向同性子空间}[totally isotropic subspace]。
\end{definition}

各向同性只要求一个非零向量的自身取值 \(q(v)=b(v,v)\) 为零；属于根空间则要求 \(b(v,w)=0\) 对所有 \(w\in V\) 成立。因此根空间中的非零向量必各向同性，反过来却不成立。全各向同性要求子空间中每个向量的二次取值为零，由极化公式，这等价于 \(b|_{T\times T}=0\)，比仅指定若干基向量各自取值为零更强。非退化空间的根空间为零，但仍可以有非零全各向同性子空间，例如上一节的 \(x^2-y^2\)。

\begin{definition}\label{b2:def:hyperbolic-plane}
设 \(k\) 为特征不为二的域，\((V,q)\) 为二维 \(k\) 二次空间，\(b=\beta_q/2\)。如果 \(V\) 有基 \(e,f\)，满足 \(q(e)=q(f)=0\)、\(b(e,f)=1\)，就称 \((V,q)\) 为\term[shuangqu-pingmian]{双曲平面}[hyperbolic plane]，记为 \(H\)。它的表达式及矩阵为
\[
q(xe+yf)=2xy,\qquad B_H=\begin{pmatrix}0&1\\1&0\end{pmatrix}.
\]
若一个 \(k\) 二次空间等距于若干个 \(H\) 的正交直和，就称其为\term[shuangqu-kongjian]{双曲空间}[hyperbolic space]，并记 \(H^{\perp r}\) 为 \(r\) 个双曲平面的正交直和。
\end{definition}
\symidx{\(H\)、\(H^{\perp r}\)：双曲平面及其正交直和}

表达式中的因子二来自本书的约定 \(q(v)=b(v,v)\)：两个混合项分别为 \(xyb(e,f)\) 与 \(xyb(f,e)\)，各等于 \(xy\)。若把未归一化的极化型 \(\beta_q\) 在基向量间的值取为一，才会得到 \(xy\) 的表达式。

\ProseParagraphBreak
所有这样给出的平面都等距，因为把指定的 \(e,f\) 相互对应便保持矩阵。\(\det B_H=-1\ne0\)，所以它非退化，却包含全各向同性直线 \(ke\) 与 \(kf\)。这两条直线的和不是全各向同性子空间，因为 \(q(e+f)=2\ne0\)；各向同性向量 \(e\) 也不属于根空间，因为 \(b(e,f)=1\)。基 \(e+f/2,e-f/2\) 相互正交，二次取值为一和负一，因而
\[
H\cong\langle1,-1\rangle.
\]
若 \(k=\mathbb R\)，这正是一个正方向与一个负方向组成的部分。

\begin{lemma}[双曲平面的分离]\label{b2:lem:hyperbolic-plane-split}
设 \(k\) 为特征不为二的域，\((V,q)\) 为有限维非退化 \(k\) 二次空间，\(b=\beta_q/2\)。若 \(e\in V\) 满足 \(e\ne0\) 且 \(q(e)=0\)，则存在 \(f\in V\)，使 \(e,f\) 生成一个双曲平面 \(H\)，并有 \(V=H\perp H^\perp\)，其中正交补非退化。
\end{lemma}
\begin{proof}
由于 \(q(e)=0\)，直线 \(ke\) 上的限制型是零型，不能把它单独取作非退化正交分量，也不能像分离非零对角项时那样除以 \(q(e)\)。我们要给它补上另一个各向同性方向，使两个方向之间的配对非零。

非退化性保证某个 \(f_0\) 满足 \(b(e,f_0)\ne0\)。缩放以后令此值为一。先尝试 \(f=f_0+te\)，其中 \(t\in k\)。由 \(q(v)=b(v,v)\)、\(q(e)=0\) 及 \(b(e,f_0)=b(f_0,e)=1\)，有
\[
b(e,f)=1,\qquad
q(f)=q(f_0)+2t\,b(f_0,e)+t^2q(e)=q(f_0)+2t.
\]
所以保持 \(b(e,f)=1\) 的同时，让 \(q(f)=0\) 只需求解线性方程 \(q(f_0)+2t=0\)。因为特征不为二，取 \(t=-q(f_0)/2\)，即
\[
f=f_0-\frac{q(f_0)}2e.
\]
两向量的配对矩阵为 \(B_H\)，所以它们线性无关且限制型非退化。现在应用\cref{b2:thm:orthogonal-complement}，即得所需正交分解与正交补的非退化性。
\end{proof}

\begin{proposition}\label{b2:prop:isotropic-subspace-hyperbolic}
设 \(k\) 为特征不为二的域，\((V,q)\) 为有限维非退化 \(k\) 二次空间。若 \(T\subseteq V\) 为 \(d\) 维全各向同性子空间，则 \(T\) 包含于某个等距于 \(H^{\perp d}\) 的非退化子空间。
\end{proposition}
\begin{proof}
对 \(d\) 归纳，\(d=0\) 时取零空间。\(d>0\) 时，取 \(0\ne e\in T\)，由\cref{b2:lem:hyperbolic-plane-split}补上 \(f\)，令 \(H=ke+kf\)。因为 \(T\) 全各向同性，其中每个 \(t\) 已与 \(e\) 正交；要让它落入 \(H^\perp\)，还须消去它与 \(f\) 的配对。利用 \(b(e,f)=1\)，沿 \(e\) 调整，置
\[
t'=t-b(t,f)e.
\]
此时 \(b(t',e)=0\)，且 \(b(t',f)=b(t,f)-b(t,f)b(e,f)=0\)。又 \(t'\in T\)，所以每个 \(t\) 都有表示 \(t=b(t,f)e+t'\)。由于 \(e\notin H^\perp\)，这个表示给出直和
\[
T=ke\oplus\bigl(T\cap H^\perp\bigr).
\]
交的维数为 \(d-1\)，仍全各向同性；在非退化的 \(H^\perp\) 内应用归纳假设，再加上 \(H\)，便得到所求子空间。
\end{proof}

\subsection{Witt 分解与消去定理}
\label{b2:subsec:1003:02}

不断分离双曲平面，每次维数减少二，便能得到分解的存在性。要说明所得各向异性部分不依赖选择，需要先证明正交直和可以消去共同的非退化部分。这里的困难是，\(U\perp V\) 到 \(U\perp W\) 的等距同构未必把指定的 \(U\) 分量送到另一边的 \(U\) 分量。将 \(U\) 对角化后，只需逐条对齐共同的非退化直线，再把映射限制到正交补；下面的反射正是用来调整这些直线的像。

\begin{lemma}[反射与等长向量]\label{b2:lem:orthogonal-reflection}
设 \(k\) 为特征不为二的域，\((V,q)\) 为有限维 \(k\) 二次空间，\(b=\beta_q/2\)。若 \(z\in V\) 满足 \(q(z)\ne0\)，则 \(V=kz\perp(kz)^\perp\)。在 \(kz\) 上取负恒等、在 \((kz)^\perp\) 上取恒等所定义的映射 \(\tau_z\) 是等距自同构，称为沿 \(z\) 的\term[zhengjiao-fanshe]{反射}[reflection]。它的公式是
\[
\tau_z(v)=v-\frac{2b(v,z)}{q(z)}z.
\]
若 \(u,v\in V\) 满足 \(q(u)=q(v)=a\ne0\)，则存在至多两个这样的反射的复合，将 \(u\) 送到 \(v\)。
\end{lemma}
\begin{proof}
因为 \(q(z)\ne0\)，直线 \(kz\) 的限制型非退化，\cref{b2:thm:orthogonal-complement}给出 \(V=kz\perp(kz)^\perp\)，这里不要求整个 \(V\) 非退化。对任意 \(v\in V\)，这一分解具体写成
\[
v=\frac{b(v,z)}{q(z)}z+v_\perp,
\qquad v_\perp\in(kz)^\perp.
\]
只把第一项取负，得到 \(-\dfrac{b(v,z)}{q(z)}z+v_\perp=v-\dfrac{2b(v,z)}{q(z)}z\)，这就导出了所给公式。两个正交分量上的作用均保持二次取值，且平方为恒等，所以 \(\tau_z\) 是等距自同构。

对等长向量，计算得
\[
q(u-v)=2\bigl(a-b(u,v)\bigr),\qquad
q(u+v)=2\bigl(a+b(u,v)\bigr).
\]
两者之和为 \(4a\ne0\)，故至少一个非零。若 \(q(u-v)\ne0\)，则
\[
\frac{2b(u,u-v)}{q(u-v)}=1,
\qquad \tau_{u-v}(u)=v.
\]
若 \(q(u-v)=0\)，则 \(q(u+v)\ne0\)，同样计算得 \(\tau_{u+v}(u)=-v\)。再施加 \(\tau_v\)，便把 \(-v\) 送到 \(v\)。这里 \(q(v)=a\ne0\)，所以第二个反射确有定义。\(u=v\) 时也可直接取恒等映射。
\end{proof}

这项反射构造可与 Roman\cite[Chapter 11, Theorem 11.32]{Roman2008AdvancedLinearAlgebra}对照。非零长度的条件十分具体：它既保证至少一个 \(u\pm v\) 能定义反射，也保证必要时能再沿 \(v\) 反射。各向同性向量不能直接使用这个论证。

\begin{theorem}[Witt 消去]\label{b2:thm:witt-cancellation}
设 \(k\) 为特征不为二的域，\(U,V,W\) 为有限维非退化 \(k\) 二次空间。若
\[
U\perp V\cong U\perp W,
\]
则 \(V\cong W\)。
\end{theorem}
\begin{proof}
先设共同部分一维。在两边分别取基向量 \(e,f\)，使 \(q(e)=q(f)=a\ne0\)，并给定等距同构
\[
\phi:ke\perp V\longrightarrow kf\perp W.
\]
目标空间中的 \(\phi(e)\) 与 \(f\) 等长且长度非零。由\cref{b2:lem:orthogonal-reflection}，有目标空间的等距自同构 \(\rho\)，使 \(\rho\phi(e)=f\)。因此 \(\rho\phi\) 将 \((ke)^\perp=V\) 送入 \((kf)^\perp=W\)：对 \(v\in V\)，保持配对给出 \(b\bigl(\rho\phi(v),f\bigr)=b(v,e)=0\)。反向应用它的逆同构，得到像恰为 \(W\)，故限制映射是 \(V\to W\) 的等距同构。

一般情形中，由\cref{b2:thm:quadratic-diagonalization}，共同非退化部分可写成
\[
U=\langle a_1\rangle\perp\cdots\perp\langle a_d\rangle,
\qquad a_i\ne0.
\]
依次对两边共同的一维部分应用已证情形。每次余下的空间仍是非退化正交直和，故下一次消去仍满足假设。经过 \(d\) 次即得 \(V\cong W\)；\(d=0\) 时本来就是该结论。
\end{proof}

上述结果称为\term[Witt-xiaoqu-dingli]{Witt 消去定理}[Witt cancellation theorem]。

\begin{theorem}[Witt 分解]\label{b2:thm:witt-decomposition}
设 \(k\) 为特征不为二的域，\((V,q)\) 为有限维非退化 \(k\) 二次空间，则 \(V\) 有分解
\[
V\cong H^{\perp r}\perp A,
\]
其中 \(A\) 各向异性。整数 \(r\) 与 \(A\) 的等距类型唯一。\(r\) 等于全各向同性子空间的最大维数，称为\term[Witt-zhishu]{Witt 指数}[Witt index]。
\end{theorem}
\begin{proof}
若有非零各向同性向量，便用\cref{b2:lem:hyperbolic-plane-split}取出一个双曲平面。在非退化正交补上重复，每次维数减少二，因此有限步后剩余部分 \(A\) 没有非零各向同性向量。若原空间已经各向异性，则取 \(r=0\)。

设还有 \(V\cong H^{\perp s}\perp B\)，其中 \(B\) 各向异性。不妨 \(r\ge s\)。由\cref{b2:thm:witt-cancellation}消去 \(s\) 个双曲平面，得到 \(H^{\perp(r-s)}\perp A\cong B\)。若 \(r>s\)，左侧含非零各向同性向量，右侧却没有，矛盾。因此 \(r=s\)，继续用所得等距即得 \(A\cong B\)。

所列双曲平面中的全部 \(e_i\) 线性无关，且 \(q(e_i)=0\)、不同平面的向量之间配对为零，所以它们生成一个 \(r\) 维全各向同性子空间。反过来，任意 \(d\) 维全各向同性子空间由\cref{b2:prop:isotropic-subspace-hyperbolic}包含在某个 \(H^{\perp d}\) 中。将其非退化正交补再作已经证明存在的 Witt 分解，得到整个 \(V\) 的一个分解，其中双曲平面的数目至少为 \(d\)。由唯一性，它恰为 \(r\)，故 \(d\le r\)。
\end{proof}

这一分解称为\term[Witt-fenjie]{Witt 分解}[Witt decomposition]。若 \(q\) 退化，记 \(R=\operatorname{rad}b\)。对 \(z\in R\)，有 \(q(z)=b(z,z)=0\) 及 \(b(v,z)=0\)，所以 \(q(v+z)=q(v)\)。因此商空间上有良定义的二次型
\[
\bar q(v+R)=q(v).
\]
它非退化：若 \(v+R\) 与全部商向量正交，就有 \(b(v,w)=0\) 对所有 \(w\in V\) 成立，从而 \(v\in R\)、\(v+R=0\)。根空间 \(R\)、商空间 \(V/R\) 及这个商型都由原型直接确定，不需要选取补空间。

\ProseParagraphBreak
要在 \(V\) 内写出直和分解，可取 \(R\) 的一个普通线性补空间 \(V_0\)。自然映射 \(V_0\to V/R\) 是线性同构，且保持二次取值，所以 \(V_0\) 的限制型非退化。对它应用 Witt 分解，得到
\[
V=R\perp V_0
\cong R\perp H^{\perp r}\perp A.
\]
第一项是零型。与第一章的商空间和补空间之别相同，\(V_0\) 在 \(V\) 中的具体位置一般不唯一，但每个这样的补空间都等距于同一个商型；商空间的 Witt 指数与各向异性部分的等距类型因而由原二次型确定。

\subsection{Witt 群的概念}
\label{b2:subsec:1003:03}

正交直和给出的等距类可以相加，但对非零 \(V\)，\((V,q)\perp(V,-q)\) 的底层空间是 \(V\oplus V\)，维数仍为 \(2\dim V\)，不能等距于零空间。一维情形已经揭示了应忽略什么：对 \(a\in k^\times\)，在 \(\langle a,-a\rangle\) 中取
\[
e=(1,1),\qquad f=\left(\frac1{2a},-\frac1{2a}\right).
\]
直接计算得 \(q(e)=q(f)=0\)、\(b(e,f)=1\)。这两个向量线性无关，因为把 \(ce+df=0\) 分别与 \(e,f\) 配对便得 \(d=c=0\)。所以 \(\langle a,-a\rangle\cong H\)。若允许增删双曲平面而不改变类，这对互为负号的一维型之和便为零；对角化后同样的配对将使 \((V,-q)\) 成为 \((V,q)\) 的加法逆元。

\begin{definition}\label{b2:def:witt-equivalence}
设 \(k\) 为特征不为二的域，\(V,W\) 为有限维非退化 \(k\) 二次空间。如果存在整数 \(r,s\ge0\)，使
\[
V\perp H^{\perp r}\cong W\perp H^{\perp s},
\]
就称 \(V\) 与 \(W\)\term[Witt-dengjia]{Witt 等价}[Witt equivalent]。所有等价类组成的集合记为 \(W(k)\)。在类上定义 \([V]+[W]=[V\perp W]\)。
\end{definition}

\begin{theorem}\label{b2:thm:witt-group}
设 \(k\) 为特征不为二的域，\(W(k)\) 为有限维非退化 \(k\) 二次空间按添加双曲平面所得的 Witt 等价类集合，则 Witt 等价是等价关系，正交直和给出 \(W(k)\) 上良定义的交换群运算。零元为零空间的类，\((V,q)\) 的负元由 \((V,-q)\) 表示。每个类恰有一个各向异性代表的等距类型。
\end{theorem}
\begin{proof}
自反性和对称性直接成立。若 \(V\perp H^{\perp r}\cong W\perp H^{\perp s}\)，又 \(W\perp H^{\perp u}\cong X\perp H^{\perp v}\)，则给第一个等距加上 \(H^{\perp u}\)，给第二个加上 \(H^{\perp s}\)，得到
\[
V\perp H^{\perp(r+u)}\cong X\perp H^{\perp(v+s)}.
\]
故关系传递。对两个等距分别取正交直和，便证明加法与代表选择无关。直和的结合、交换及零空间的单位性质给出相应的群公理。

还须构造负元。将 \(q\) 对角化为 \(\langle a_1,\ldots,a_n\rangle\)，其中 \(a_i\ne0\)。\(q\perp(-q)\) 重排分量后为若干个 \(\langle a_i,-a_i\rangle\)，由定义前的一维计算，每个分量都等距于 \(H\)。因此 \((V,q)\perp(V,-q)\cong H^{\perp n}\)，其类为零，给出负元。

\cref{b2:thm:witt-decomposition}说明每个类都由各向异性部分表示。若各向异性 \(A,B\) Witt 等价，则 \(A\perp H^{\perp r}\cong B\perp H^{\perp s}\) 的两边已是 Witt 分解，唯一性迫使 \(r=s\) 且 \(A\cong B\)。
\end{proof}

这个群称为域 \(k\) 的\term[Witt-qun]{Witt 群}[Witt group]。记号中的集合没有大小困难：每个有限维型都可以由某个有限对称矩阵表示，全部这些矩阵组成集合，再取等价类即可。
\symidx{\(W(k)\)：特征不为二的域的 Witt 群}

\begin{proposition}\label{b2:prop:witt-real-complex}
对有限维非退化实二次空间，正惯性指数与负惯性指数之差给出群同构 \(W(\mathbb R)\cong\mathbb Z\)。对有限维非退化复二次空间，维数的奇偶性给出群同构 \(W(\mathbb C)\cong\mathbb Z/2\mathbb Z\)。
\end{proposition}
\begin{proof}
实型由\cref{b2:thm:sylvester-inertia}分为 \(p_+\) 个 \(\langle1\rangle\) 与 \(p_-\) 个 \(\langle-1\rangle\)。一正一负合成双曲平面；配对后剩余系数全为一或全为负一，是各向异性型。它的维数与符号由差 \(p_+-p_-\) 决定，所以其 Witt 类只由这个差决定。符号差在正交直和下相加，加双曲平面不改变它；差为零时全部分量都能配成双曲平面，而任意整数都能由全正或全负型实现，故所得同态既单射又满射。

复数域上，先由\cref{b2:thm:quadratic-diagonalization}对角化，再利用每个非零复数都有平方根，把全部非零系数缩放为一。因此非退化型只有维数这一分类数据。\(\langle1,1\rangle\cong\langle1,-1\rangle\cong H\)，第一个等距通过一个坐标乘以 \(i\) 得到。因此偶数维型的类为零，奇数维型的类为 \(\bigl[\langle1\rangle\bigr]\)。加双曲平面不改变维数奇偶，故这个非零类不可能等于零，得到所述二阶群。
\end{proof}

张量积还能在这个群上定义乘法。\appref{b2:app:B}的\cref{b2:thm:B-witt-ring}证明该乘法对Witt类良定义，\cref{b2:prop:B-discriminant-witt}进一步讨论维数奇偶、判别式和基本理想；这些是本节加法理论之后的补充。

% !TeX root = ../../Book2.tex
% 纲要标题：### 10.4 经典群及复数结构
% 纲要位置：计划纲要.md，第 1392 行。

\section{经典群及复数结构}
\label{b2:sec:1004}

一个型给空间增加了结构，保持这种结构的可逆线性变换便组成一个群。不同的型给出不同的经典群；复数域上还须区别双线性配对与带共轭的配对。本节最后的 Clifford 代数说明二次型也能直接参与一个结合代数的定义。

% 纲要要点（供目录核对）：
%
% * 正交群、辛群及酉群
% * Hermitian 型与复数结构；用实部和乘以i恢复完整配对
% * 完整证明非退化型的伴随算子存在唯一性、反乘法性质及像与核关系；型的相似变换与乘子的实际像
% * 区别型的相似变换、矩阵共轭与一般线性群；正交、辛、Hermitian 条件分别写清
% * Clifford 代数的构造作为选读延伸；零型回到外代数，证明双曲平面及一维模型
% #### 10.4.3* Clifford 代数的构造
%
% ---
%

\subsection{正交群、辛群与酉群}
\label{b2:subsec:1004:01}

设 \(\operatorname{char}k\ne2\)，\((V,q)\) 为有限维非退化二次空间。保持 \(q\) 的线性自同构组成\term[zhengjiao-qun]{正交群}[orthogonal group]
\[
\operatorname O(V,q)=\{\,g\in\operatorname{GL}(V)\mid q(gv)=q(v)\text{ 对所有 }v\in V\,\}.
\]
由\eqref{b2:eq:normalized-polar}，保持 \(q\) 等价于保持 \(b_q\)。在 Gram 矩阵为 \(B\) 的基下，\footnote{格拉姆（Jørgen Pedersen Gram，1850-1916），丹麦数学家与精算师，研究向量内积和正交化。}这个群写为
\[
\operatorname O(B)=\{\,G\in\operatorname{GL}_n(k)\mid G^{\mathsf T}BG=B\,\}.
\]
保持非退化交替型 \(b\) 的线性自同构则组成\term[xin-qun]{辛群}[symplectic group]。在\cref{b2:thm:alternating-form-normal}给出的辛基下，记为
\[
\operatorname{Sp}_{2m}(k)=\{\,G\in\operatorname{GL}_{2m}(k)\mid G^{\mathsf T}J_mG=J_m\,\}.
\]
辛群的这一定义允许任意特征。
\symidx{\(\operatorname O(V,q)\)、\(\operatorname{Sp}_{2m}(k)\)：正交群与辛群}

这些集合确为群：恒等变换保持型；保持型的两个映射的复合也保持型；在保持型的等式中代入逆映射的像，就得到逆映射也保持型。改变基只改变群在一般线性群中的矩阵实现，不改变其抽象含义。

\ProseParagraphBreak
对 \(G^{\mathsf T}BG=B\) 取行列式，因 \(\det B\ne0\)，得到 \((\det G)^2=1\)。于是 \((\det G-1)(\det G+1)=0\)；域中没有零因子，故 \(\det G\in\{1,-1\}\)，而特征不为二保证这两个值不同。行列式为一的子群称为\term[teshu-zhengjiao-qun]{特殊正交群}[special orthogonal group]，记为 \(\operatorname{SO}(V,q)\)。当 \(V\ne0\) 时，\cref{b2:thm:quadratic-diagonalization}给出非零长度向量，沿它的反射在该直线上为负恒等、在正交补上为恒等，所以行列式为负一。故此时特殊正交群在正交群中指数为二。

\ProseParagraphBreak
复向量空间 \(\mathbb C^n\) 上的配对
\[
h_0(z,w)=\overline z^{\mathsf T}w=\sum_i\overline{z_i}w_i
\]
在第二变量线性，在第一变量共轭线性。保持它的复线性变换组成\term[you-qun]{酉群}[unitary group]
\[
\operatorname U(n)=\{\,G\in\operatorname{GL}_n(\mathbb C)\mid G^*G=I_n\,\},
\qquad G^*=\overline G^{\mathsf T}.
\]
于是 \(|\det G|^2=1\)。复正交型 \(z^{\mathsf T}w\) 对两个变量都线性，\(h_0\) 却在第一变量带有共轭。这个差别已经在乘 \(i\) 时显现：
\[
(iz)^{\mathsf T}(iw)=-z^{\mathsf T}w,\qquad
\overline{iz}^{\mathsf T}(iw)=\overline z^{\mathsf T}w.
\]
所以 \(iI_n\) 保持 \(h_0\)，但当 \(n>0\) 时不属于 \(\operatorname O(I_n)\)。
\symidx{\(\operatorname U(n)\)、\(G^*=\overline G^{\mathsf T}\)：标准酉群与共轭转置}

\ProseParagraphBreak
二维辛群还有一个直接可算的描述。若 \(G=\begin{psmallmatrix}a&b\\c&d\end{psmallmatrix}\)，相乘得到
\[
G^{\mathsf T}\begin{pmatrix}0&1\\-1&0\end{pmatrix}G
=(ad-bc)\begin{pmatrix}0&1\\-1&0\end{pmatrix}.
\]
所以 \(\operatorname{Sp}_2(k)=\operatorname{SL}_2(k)\)，且这个等式在特征二中仍成立。它来自这个二维计算，不能把正交群与辛群的一般定义混为一谈。

\subsection{Hermitian 型、伴随与型的相似变换}
\label{b2:subsec:1004:02}

\begin{definition}\label{b2:def:hermitian-form}
设域 \(K\) 配有一个\term[yu-duihe]{对合}[involution]，即自同构 \(\sigma\) 满足 \(\sigma^2=\operatorname{id}\)，并记 \(\overline a=\sigma(a)\)。给定 \(K\) 向量空间 \(V\)，如果双可加映射 \(h:V\times V\to K\) 满足
\[
h(au,bv)=\overline a\,b\cdot h(u,v)
\]
且 \(h(v,u)=\overline{h(u,v)}\) 对所有 \(a,b\in K\)、\(u,v\in V\) 成立，就称 \(h\) 为 \(V\) 上的\term[Hermitian-xing]{Hermitian 型}[Hermitian form]（也称厄米型）。
\termidx[emi-xing]{厄米型}
若 \(h(v,w)=0\) 对全部 \(w\) 成立只能有 \(v=0\)，称 \(h\) 非退化。保持 \(h\) 的 \(K\) 线性自同构组成的群记为 \(\operatorname U(V,h)\)。
\end{definition}

恒等对合给出对称双线性型；复共轭给出通常的 Hermitian 型。本书固定第一变量共轭线性、第二变量线性的约定。若 \(H=\bigl(h(e_i,e_j)\bigr)\)，则
\[
h(u,v)=\overline x^{\mathsf T}Hy,\quad H^*=H,
\qquad H'=P^*HP.
\]
最后一个公式由 \(x=Px',y=Py'\) 直接代入得到。对角元都属于固定域 \(K^\sigma=\{\,a\in K\mid\sigma(a)=a\,\}\)：恒等对合的固定域为 \(K\)，复共轭的固定域则为 \(\mathbb R\)。保持型的矩阵满足 \(G^*HG=H\)。

\ProseParagraphBreak
分离一维正交分量需要先找到 \(h(v,v)\ne0\) 的向量。若对合恒等且特征为二，非零型却可能是交替型；例如 \(K^2\) 上的 \(h(u,v)=u_1v_2+u_2v_1\) 非零，但每个 \(h(v,v)\) 都为零。这样的型须使用辛型的二维块，不能从非零长度的直线开始分离。这正是下面假设所排除的情形。

\begin{proposition}\label{b2:prop:hermitian-diagonalization}
设 \(K\) 为配有对合 \(\sigma\) 的域，\(K^\sigma=\{a\in K\mid\sigma(a)=a\}\)，\(V\) 为有限维 \(K\) 向量空间，\(h:V\times V\to K\) 为 Hermitian 型。若 \(\sigma\ne\operatorname{id}\)，或 \(\operatorname{char}K\ne2\)，则 \(h\) 有正交基，且对角系数属于 \(K^\sigma\)。在 \(K=\mathbb C\)、\(\sigma\) 为复共轭的情形，对角系数可以归一化为一、负一与零，其个数唯一。
\end{proposition}
\begin{proof}
先证明非零的型必有 \(h(v,v)\ne0\) 的向量。反设全部对角取值为零。令 \(c=h(u,v)\)，从 \(h(u+v,u+v)=0\) 得 \(c+\overline c=0\)。再对 \(u+av\) 作同样展开，得到
\[
ac+\overline a\,\overline c=(a-\overline a)c=0.
\]
若对合非恒等，选择 \(a\ne\overline a\) 便得 \(c=0\)。若对合恒等而特征不为二，第一式就是 \(2c=0\)，也得 \(c=0\)。任意 \(u,v\) 的配对均为零，与型非零矛盾。

若型为零，任意基均正交。否则取 \(e\) 使 \(d=h(e,e)\ne0\)。对每个 \(v\)，向量 \(v-\dfrac{h(e,v)}d\,e\) 与 \(e\) 正交，因为第二变量线性给出 \(h(e,v)-h(e,v)d/d=0\)；由 Hermitian 对称性，两种次序的配对都为零。又 \(Ke\) 与其正交补的交为零，故分离出这个非退化一维分量，再按维数归纳。对角元被 \(\sigma\) 固定，已经由 Hermitian 对称性保证。

在复数情形，这些系数均为实数，缩放基向量可把非零系数改为正负一。若正项空间维数为 \(p_+\)，任意正定复子空间投影到它都单射，因为投影为零的向量的二次取值非正。故 \(p_+\) 是正定子空间最大复维数。负项的个数对 \(-h\) 作同样论证确定，零项数由总维数减去二者得到。这与\cref{b2:thm:sylvester-inertia}中的论证相同，只把实坐标平方换成复坐标模平方。
\end{proof}

非退化复 Hermitian 型的惯性为 \((p_+,p_-,0)\) 时，对应群常记为 \(\operatorname U(p_+,p_-)\)，其标准矩阵为 \(\operatorname{diag}(I_{p_+},-I_{p_-})\)。标准酉群是 \(p_-=0\) 的情形。这里的“酉”只表示保持指定 Hermitian 型，不自动表示该型正定。标准酉群还可以完全用实线性变换描述：忘去 \(\mathbb C^n\) 的复标量后，保留实配对 \(\operatorname{Re}h_0\) 与乘 \(i\) 的实线性算子，就能恢复原来的复结构和配对。

\begin{proposition}\label{b2:prop:unitary-real-complex-structure}
设 \(n\ge0\)，\(h_0(z,w)=\sum_{j=1}^n\overline{z_j}w_j\) 为 \(\mathbb C^n\) 上的标准 Hermitian 型。将 \(\mathbb C^n\) 看成实空间，记 \(Jz=iz\)、\(g(z,w)=\operatorname{Re}h_0(z,w)\)，则 \(J^2=-I\)，而保持 \(h_0\) 的复线性自同构群 \(\operatorname U(n)\) 恰为保持 \(g\) 并与 \(J\) 交换的实线性自同构组成的群。
\end{proposition}
\begin{proof}
实线性映射与 \(J\) 交换，当且仅当它还保持乘 \(i\)，也就当且仅当它复线性。又由第一变量的共轭线性，\(h_0(Jz,w)=-i h_0(z,w)\)，所以
\[
g(Jz,w)=\operatorname{Im}h_0(z,w),\qquad
h_0(z,w)=g(z,w)+i\,g(Jz,w).
\]
若映射保持 \(h_0\)，它显然保持实部，且作为复线性映射与 \(J\) 交换。反过来，若 \(A\) 保持 \(g\) 并与 \(J\) 交换，则
\[
g(JAz,Aw)=g(AJz,Aw)=g(Jz,w),
\]
所以它也保持 \(h_0\) 的虚部，连同实部便保持 \(h_0\)。
\end{proof}

若实向量空间 \(W\) 上的实线性算子 \(J:W\to W\) 满足 \(J^2=-I_W\)，就称 \(J\) 为 \(W\) 上的一个\term[fushu-jiegou]{复数结构}[complex structure]：规定 \((a+bi)v=av+bJv\)，就使原实空间成为复向量空间。结合律由 \(J^2=-I_W\) 验证。若 \(W\) 的实维数有限，复维数也有限；一组复基 \(w_1,\ldots,w_r\) 对应实基 \(w_1,Jw_1,\ldots,w_r,Jw_r\)，所以 \(\dim_{\mathbb R}W=2\dim_{\mathbb C}W\) 必为偶数。上面的命题因而把酉群同时解释为一种实正交群的子群和一种复线性群。

\ProseParagraphBreak
型还可以用来给线性算子定义伴随。下面把对称型、交替型和 Hermitian 型一并处理；符号中的共轭在前两种情形取恒等，在第三种情形取指定的域对合。

\begin{proposition}\label{b2:prop:form-adjoint-operator}
设 \(K\) 为配有对合 \(\sigma:a\mapsto\bar a\) 的域，\(V\) 为有限维 \(K\) 向量空间，\(\varepsilon\in\{1,-1\}\)。设 \(b:V\times V\to K\) 为非退化的 \(\varepsilon\)-Hermitian 型，即 \(b\) 双可加，第一变量共轭线性、第二变量线性，并满足
\[
b(au,cv)=\bar a\,c\cdot b(u,v),\qquad
b(v,u)=\varepsilon\overline{b(u,v)}
\]
对所有 \(a,c\in K\)、\(u,v\in V\) 成立，且 \(b(u,v)=0\) 对所有 \(v\) 成立只能有 \(u=0\)。对每个 \(K\) 线性算子 \(T:V\rightarrow V\)，存在唯一的 \(K\) 线性算子 \(T^\dagger:V\rightarrow V\) 满足
\[
b(Tu,v)=b(u,T^\dagger v)\qquad(u,v\in V).
\]
称 \(T^\dagger\) 为 \(T\) 关于 \(b\) 的\term[bansui-suanzi]{伴随算子}[adjoint operator]。选定基，令 \(H=(b(e_i,e_j))\)，并以相同字母表示算子的矩阵，则 \(H\) 可逆、\(H^*=\varepsilon H\)，其中 \(X^*=\overline X^{\mathsf T}\)。对任意另一 \(K\) 线性算子 \(S:V\rightarrow V\)，有
\[
T^\dagger=H^{-1}T^*H,\qquad
(ST)^\dagger=T^\dagger S^\dagger,\qquad
(T^\dagger)^\dagger=T.
\]
伴随运算可加，并满足 \((aT)^\dagger=\bar aT^\dagger\)。对任意子空间 \(W\le V\)，记 \(W^\perp=\{\,v\in V\mid b(w,v)=0\text{ 对所有 }w\in W\,\}\)，则
\[
(\operatorname{im}T)^\perp=\ker T^\dagger,\qquad
\operatorname{rank}T^\dagger=\operatorname{rank}T.
\]
\end{proposition}
\begin{proof}
固定第二变量时，\(u\mapsto b(u,v)\) 是共轭线性泛函，不能在非恒等对合下把它当作普通 \(V^*\) 中的元素。令
\[
C_\sigma(V,K)=
\{\,\ell:V\to K\mid \ell\text{ 可加，且 }\ell(au)=\bar a\,\ell(u)
\ (a\in K,\ u\in V)\,\}
\]
按逐点加法与标量乘法组成 \(K\) 向量空间。选任意一组 \(V\) 的基，共轭线性泛函由它在基上的值唯一确定，这些值又可任意指定，所以 \(\dim_K C_\sigma(V,K)=\dim_KV\)。映射
\[
\Psi:V\longrightarrow C_\sigma(V,K),\qquad
w\longmapsto \bigl(u\longmapsto b(u,w)\bigr)
\]
由第二变量线性而 \(K\) 线性。若 \(\Psi(w)=0\)，则 \(b(u,w)=0\) 对所有 \(u\) 成立，由 \(\varepsilon\)-Hermitian 对称性也有 \(b(w,u)=0\)，非退化性给出 \(w=0\)。有限维且等维使 \(\Psi\) 成为同构。

对每个 \(v\in V\)，\(\ell_v(u)=b(Tu,v)\) 属于 \(C_\sigma(V,K)\)，故存在唯一 \(w\) 使 \(\Psi(w)=\ell_v\)。定义 \(T^\dagger v=w\)，便得到要求的配对等式。又 \(v\mapsto\ell_v\) 是 \(K\) 线性，后接 \(\Psi^{-1}\) 说明 \(T^\dagger\) 线性；每个 \(v\) 所对应的 \(w\) 唯一也给出算子唯一性。这与\cref{b2:prop:finite-perfect-pairing}的有限维表示机制相同，只是第一变量的泛函改为共轭线性泛函。

选定基以后，要求的配对等式成为
\[
\overline u^{\mathsf T}T^*Hv
=\overline u^{\mathsf T}HT^\dagger v.
\]
对全部基向量 \(u,v\) 成立，当且仅当 \(T^*H=HT^\dagger\)。非退化性使 \(H\) 可逆，故得到矩阵公式。复合反序也可直接从配对看出：
\[
b(STu,v)=b(Tu,S^\dagger v)
=b(u,T^\dagger S^\dagger v).
\]
由唯一性，\((ST)^\dagger=T^\dagger S^\dagger\)。矩阵中的共轭转置也反转乘法，与这个内在公式一致。再由 \(H^*=\varepsilon H\) 以及 \(\varepsilon^2=1\)，有
\[
(T^\dagger)^\dagger
=H^{-1}H^*T(H^*)^{-1}H=T.
\]
可加性和标量公式同样直接来自共轭转置。

向量 \(v\) 属于 \((\operatorname{im}T)^\perp\)，当且仅当 \(b(Tu,v)=0\) 对所有 \(u\) 成立；伴随等式把这个条件改写为 \(b(u,T^\dagger v)=0\) 对所有 \(u\) 成立。由第二变量的非退化性，这恰好等价于 \(T^\dagger v=0\)，证明正交补与核的等式。

对任意 \(W\le V\)，限制映射 \(C_\sigma(V,K)\to C_\sigma(W,K)\) 满射：把 \(W\) 的一组基扩充为 \(V\) 的基，并在新增基向量上任意指定值，即可延伸共轭线性泛函。因此 \(\Psi\) 后接限制的映射 \(V\to C_\sigma(W,K)\) 满射，核为 \(W^\perp\)，故
\(\dim W^\perp=\dim V-\dim W\)。取 \(W=\operatorname{im}T\)，并对 \(T^\dagger\) 使用秩与核的维数公式，得到
\[
\operatorname{rank}T^\dagger
=\dim V-\dim\ker T^\dagger
=\dim\operatorname{im}T=\operatorname{rank}T.
\]
\end{proof}
\symidx{\(T^\dagger\)：相对于指定非退化型的伴随算子}

例如，实标准内积给出 \(T^\dagger=T^{\mathsf T}\)，复标准 Hermitian 型给出 \(T^\dagger=T^*\)。对辛矩阵 \(J_m\)，则有 \(T^\dagger=J_m^{-1}T^{\mathsf T}J_m\)。这都是整个线性算子的伴随，不是由代数余子式组成的伴随矩阵。这个反转乘法的运算将在\secref{b2:sec:B03}作为代数上的对合进一步讨论。

\ProseParagraphBreak
算子换基用矩阵相似 \(A\mapsto P^{-1}AP\)，型的换基则用合同 \(H\mapsto P^*HP\)；二者都是重新描述同一个对象。型的换基不要求 \(P\) 保持原矩阵 \(H\)。若把可逆矩阵 \(g\) 看作固定空间上的变换，等距条件才是 \(g^*Hg=H\)；以下的型相似变换则只要求 \(g^*Hg=cH\)，即把全部配对值同时乘以一个非零标量。它与算子的矩阵相似是不同概念。

\begin{definition}\label{b2:def:form-similitude}
设 \(K\) 为带域对合 \(a\mapsto\bar a\) 的域，固定域为 \(F=\{\,a\in K\mid\bar a=a\,\}\)，\(V=K^n\)、\(n\ge1\)。取可逆矩阵 \(H\) 满足 \(H^*=\varepsilon H\)，其中 \(X^*=\overline X^{\mathsf T}\)、\(\varepsilon\in\{1,-1\}\)，并令 \(b(u,v)=\overline u^{\mathsf T}Hv\)。若 \(g\in\operatorname{GL}_K(V)\) 满足
\[
b(gu,gv)=c\cdot b(u,v)\qquad(u,v\in V)
\]
对某个 \(c\in F^\times\) 成立，就称 \(g\) 为一个\term[xing-xiangsi-bianhuan]{型的相似变换}[similitude]，称 \(c\) 为其\term[xiangsi-chengzi]{乘子}[multiplier]，记为 \(\mu(g)\)。全部这类变换组成相似变换群。共轭为恒等且特征不为二时，对称型的相似变换群记为 \(\operatorname{GO}(b)\)；共轭为恒等时，交替型的相似变换群记为 \(\operatorname{GSp}(b)\)；对给定非平凡域对合，Hermitian 型的相似变换群记为 \(\operatorname{GU}(b)\)。特征二中的二次型相似条件须直接用二次型定义，不能只由其极化型判断。
\end{definition}

型非退化且 \(n\ge1\)，保证可以取 \(u,v\) 使 \(z=b(u,v)\ne0\)。若 \(c,c'\) 都满足相似条件，则 \(cz=c'z\)，在域中消去非零的 \(z\)，得到 \(c=c'\)，所以乘子唯一。

\ProseParagraphBreak
即使先只要求 \(c\in K^\times\)，交换这两个变量也会迫使它属于固定域。一方面，
\[
b(gv,gu)=c\,b(v,u)=c\varepsilon\bar z;
\]
另一方面，由型的对称性，
\[
b(gv,gu)=\varepsilon\overline{b(gu,gv)}
=\varepsilon\bar c\,\bar z.
\]
因 \(\varepsilon\bar z\ne0\)，这两个等式给出 \(c=\bar c\)。

\begin{proposition}\label{b2:prop:similitude-multiplier}
设 \(K\) 为带域对合 \(a\mapsto\bar a\) 的域，\(F\) 为其固定域，\(V=K^n\)、\(n\ge1\)，且 \(b(u,v)=\overline u^{\mathsf T}Hv\)，其中 \(H\) 可逆、\(H^*=\varepsilon H\)、\(\varepsilon\in\{1,-1\}\)，星号表示共轭转置。型 \(b\) 的相似变换组成群，乘子为到 \(F^\times\) 的群同态，其核恰为保持 \(b\) 的等距群。记 \(g^\dagger=H^{-1}g^*H\)，则对 \(g\in\operatorname{GL}_n(K)\)、\(c\in F^\times\)，有
\[
g\text{ 是乘子为 }c\text{ 的相似变换}
\quad\Longleftrightarrow\quad
g^*Hg=cH
\quad\Longleftrightarrow\quad
g^\dagger g=cI_n.
\]
\end{proposition}
\begin{proof}
两个相似变换依次作用时，配对值乘以两个乘子的积；恒等变换的乘子为一。把 \(u,v\) 换成 \(g^{-1}u,g^{-1}v\)，可知逆变换的乘子为 \(\mu(g)^{-1}\)。这证明群结构及同态性。乘子为一恰好是等距条件。把配对写成矩阵，再使用 \(g^\dagger=H^{-1}g^*H\)，便得两个等价式。
\end{proof}

乘子的核已经确定，接下来还需计算它的实际像，而这个像受型的性质限制。在非零空间上的正定实型 \(b(u,v)=u^{\mathsf T}v\) 上，非零 \(v\) 给出 \(\mu(g)=b(gv,gv)/b(v,v)>0\)，而 \(\sqrt c I_n\) 实现每个 \(c>0\)，故像恰为 \(\mathbb R_{>0}\)。非零空间上的复正定 Hermitian 型也有相同的乘子像。另一方面，对标准辛型取
\[
g_t=\operatorname{diag}(tI_m,I_m),\qquad t\in K^\times.
\]
直接按块相乘，得到 \(g_t^{\mathsf T}J_mg_t=tJ_m\)，所以当 \(m\ge1\) 时，辛相似变换的乘子取遍 \(K^\times\)。因此从核得到的正合序列应以乘子的实际像结束，不能无条件把它写成到全部非零标量的满射。

\OptionalSubsection{Clifford 代数的构造}{b2:subsec:1004:03}

下面恢复任意域 \(k\)，并令 \(q\) 按\cref{b2:def:quadratic-form}定义。张量代数保留向量空间的线性关系而不另加乘法关系；现在要求向量 \(v\) 的像满足平方等式 \(v^2=q(v)1\)，就是把二次型的取值直接写进乘法。将这些等式的差生成双边理想，便使它们在左右乘法下始终成立。构造允许型退化，也允许特征二；此时仍使用未经除以二的 \(\beta_q\)。

\begin{definition}\label{b2:def:clifford-algebra}
设 \(k\) 为任意域，\(V\) 为 \(k\) 向量空间，\(q:V\to k\) 为二次型。取张量代数 \(T_k(V)\)，令 \(I_q\) 为全部 \(v\otimes v-q(v)1\)、\(v\in V\) 生成的双边理想。商代数
\[
\operatorname{Cl}(V,q)=T_k(V)/I_q
\]
称为 \(q\) 的\term[Clifford-daishu]{Clifford 代数}[Clifford algebra]。从 \(V\) 到这个商代数的自然线性映射记为 \(\iota\)。
\end{definition}
\symidx{\(\operatorname{Cl}(V,q)\)：二次型的 Clifford 代数}

若 \(q=0\)，关系理想恰由所有 \(v\otimes v\) 生成。由\cref{b2:prop:exterior-algebra-universal}的张量商描述，得到自然代数同构
\[
\operatorname{Cl}(V,0)\cong\bigwedge_k V.
\]
因此外代数就是所有一次元素平方为零的 Clifford 代数；一般二次型把这些平方从零改为相应的标量。

\begin{theorem}[Clifford 代数的泛性质]\label{b2:thm:clifford-universal}
设 \(k\) 为任意域，\(V\) 为 \(k\) 向量空间，\(q:V\to k\) 为二次型，\(\iota:V\to\operatorname{Cl}(V,q)\) 为自然映射。对任意结合含幺 \(k\) 代数 \(A\)，若 \(k\) 线性映射 \(f:V\to A\) 满足 \(f(v)^2=q(v)1_A\) 对每个 \(v\in V\) 成立，则存在唯一含幺 \(k\) 代数同态
\[
\widetilde f:\operatorname{Cl}(V,q)\longrightarrow A,
\qquad \widetilde f\iota=f.
\]
此外在 Clifford 代数中有
\begin{equation}\label{b2:eq:clifford-polar-relation}
\iota(u)\iota(v)+\iota(v)\iota(u)=\beta_q(u,v)1.
\end{equation}
\end{theorem}
\begin{proof}
\cref{b2:thm:tensor-algebra-universal}将 \(f\) 唯一延伸为 \(T_k(V)\to A\) 的含幺代数同态。题设等式使 \(I_q\) 的每个生成元映到零，双边理想的全部元素也就映到零，故它通过商代数分解。反过来，任何这样的分解都给出原来的张量代数同态，所以唯一。

将 \(\iota(u+v)^2=q(u+v)1\) 展开，减去 \(\iota(u)^2=q(u)1\) 与 \(\iota(v)^2=q(v)1\)，便得到\eqref{b2:eq:clifford-polar-relation}。特征不为二时，右侧为 \(2b_q(u,v)1\)。
\end{proof}

关系 \(\iota(v)^2=q(v)1\) 记录向量的二次取值；特征不为二时，\eqref{b2:eq:clifford-polar-relation}还说明，若 \(b_q(u,v)=0\)，则 \(\iota(u)\iota(v)=-\iota(v)\iota(u)\)。乘法同时记录了二次取值与正交条件。

\ProseParagraphBreak
上述商构造给出了自然线性映射 \(\iota\) 及泛性质。\(\iota\) 的单射性和代数的维数还需由基的结构证明；在特征不为二的有限维情形，\cref{b2:thm:B-clifford-basis}给出其单射性及维数 \(2^{\dim V}\)。下面直接计算一个低维模型。

\begin{example}\label{b2:exm:clifford-hyperbolic-plane}
设 \(k\) 为特征不为二的域，\(H=ke\oplus kf\) 上的二次型为 \(q(xe+yf)=2xy\)，则 \(\operatorname{Cl}(H,q)\cong M_2(k)\)。
\end{example}
\begin{proof}
以双曲基 \(e,f\) 的像仍记为 \(e,f\)，关系为
\[
e^2=f^2=0,\qquad ef+fe=2.
\]
先说明 \(1,e,f,ef\) 生成整个代数的向量空间。任何字若出现相邻相同字母便为零；否则字母交替。长度至少三时含 \(efe\) 或 \(fef\)，而
\[
efe=e(2-ef)=2e,\qquad fef=(2-ef)f=2f.
\]
递减字长后只剩长度至多二的字，再用 \(fe=2-ef\) 即得四个所列元素。

为找到满足这些关系的矩阵，先取矩阵单位 \(E_{12},E_{21}\)。它们的平方均为零，且 \(E_{12}E_{21}+E_{21}E_{12}=I_2\)；把第一个矩阵乘二，就能实现右端的 \(2I_2\)。因此取
\[
E=\begin{pmatrix}0&2\\0&0\end{pmatrix},\qquad
F=\begin{pmatrix}0&0\\1&0\end{pmatrix}.
\]
它们满足 \(E^2=F^2=0\)、\(EF+FE=2I_2\)，因此线性映射 \(xe+yf\mapsto xE+yF\) 的平方为 \(2xyI_2\)。由\cref{b2:thm:clifford-universal}得到代数同态到 \(M_2(k)\)。\(I_2,E,F,EF\) 线性无关：两个非对角位置先确定 \(E,F\) 的系数为零，两个对角位置再确定余下系数为零，所用的二可逆。它们构成四维矩阵空间的基。来源已由四个元素生成，故同态既单射又满射。
\end{proof}

在这个例子里，二次关系产生了一个非交换矩阵代数。关系 \(v\otimes v-q(v)1\) 混合张量次数二与零，所以商代数一般不继承张量次数的自然数分次；将次数模二以后，两项却同属偶次。因而 \(I_q\) 对奇偶分次仍是齐次理想，商代数继承直和分解
\[
\operatorname{Cl}(V,q)=\operatorname{Cl}^0(V,q)\oplus
\operatorname{Cl}^1(V,q),
\qquad
\operatorname{Cl}^i\operatorname{Cl}^j\subseteq
\operatorname{Cl}^{\,i+j\ (\mathrm{mod}\ 2)}.
\]
在特征不为二的非零有限维非退化二次空间上，\cref{b2:prop:B-spin-action}再用偶部分的可逆元构造自旋群到特殊正交群的作用，并算出其核。

\begin{proposition}\label{b2:prop:clifford-line-models}
设 \(k\) 为任意域，\(V=ke\) 为一维 \(k\) 向量空间，\(a\in k\)，二次型为 \(q(xe)=ax^2\)。则存在含幺 \(k\) 代数同构
\[
\operatorname{Cl}(V,q)\cong k[t]/(t^2-a),
\]
将 \(\iota(e)\) 送到 \(t\) 的剩余类。特别地，在 \(k=\mathbb R\) 时，
\[
\operatorname{Cl}(\mathbb R e,q)\cong
\begin{cases}
\mathbb R\times\mathbb R,&a>0,\\
\mathbb C,&a<0,\\
\mathbb R[\varepsilon]/(\varepsilon^2),&a=0.
\end{cases}
\]
\end{proposition}
\begin{proof}
令 \(A=k[t]/(t^2-a)\)，\(\tau=t+(t^2-a)\)。线性映射 \(f:V\to A\)、\(xe\mapsto x\tau\) 满足
\[
f(xe)^2=x^2\tau^2=ax^2\,1_A=q(xe)1_A.
\]
由\cref{b2:thm:clifford-universal}，得到含幺代数同态
\(\alpha:\operatorname{Cl}(V,q)\to A\)，使 \(\alpha(\iota(e))=\tau\)。

反过来，一生成元的自由结合代数是 \(k[t]\)，故赋值 \(t\mapsto\iota(e)\) 给出含幺代数同态 \(k[t]\to\operatorname{Cl}(V,q)\)。关系 \(\iota(e)^2=a1\) 使 \(t^2-a\) 映到零，因而同态通过商，得到
\(\beta:A\to\operatorname{Cl}(V,q)\)。复合 \(\alpha\beta\) 固定 \(\tau\) 与标量，故在由它们生成的 \(A\) 上为恒等；复合 \(\beta\alpha\) 固定 \(\iota(e)\) 与标量，而它们生成整个 Clifford 代数，故这个复合也为恒等。这证明所述同构，对特征二同样适用。

首一二次多项式的带余除法说明，\(A\) 中每个元素唯一写成 \(x+y\tau\)。若 \(a>0\)，令 \(r=\sqrt a\)，赋值 \(\tau\mapsto(r,-r)\) 给出 \(\mathbb R\) 代数同态
\[
A\longrightarrow\mathbb R\times\mathbb R,\qquad
x+y\tau\longmapsto(x+yr,x-yr).
\]
它的逆将 \((u,v)\) 送到 \(\dfrac{u+v}2+\dfrac{u-v}{2r}\tau\)，所以是同构。若 \(a<0\)，令 \(r=\sqrt{-a}\)，赋值 \(\tau\mapsto ir\) 给出
\[
A\longrightarrow\mathbb C,\qquad x+y\tau\longmapsto x+iyr.
\]
因 \(r\ne0\)，每个复数唯一有这种表达，故也是同构。若 \(a=0\)，只需把 \(\tau\) 政名为 \(\varepsilon\)，关系就是 \(\varepsilon^2=0\)，得到双数代数。
\end{proof}

本章的 Witt 群若要连同张量积一起计算，可在读完\secref{b2:sec:1003}后转到\secref{b2:sec:B02}；型的伴随算子同代数对合的联系则见\secref{b2:sec:B03}，那一节还需要\chapref{b2:ch:17}的中心单代数与四元数。


\begin{chaptersummary}
配对使向量与泛函相联系；有限维非退化性把这种联系变成同构，继而给出正交补的维数公式。只有当子空间上的限制型非退化时，正交补才直接给出直和。型的矩阵按合同变化，算子的矩阵按相似变化，二者的分类不变量不同。

\ProseParagraphBreak
交替型由空间维数与秩共同分类。特征不为二的二次型可以对角化，但进一步的分类取决于域：实数域上，正负惯性指数与根空间维数共同分类；复数域上，空间维数与秩共同分类。一般域上，非零各向同性向量使双曲平面得以分离；反射保证共同的非退化分量可以消去，从而确定唯一的各向异性部分。Witt 群进一步忽略双曲部分，保留二次型在稳定等距下的信息。

\ProseParagraphBreak
正交群、辛群与酉群分别保持指定的型。伴随算子把保持型的条件写成 \(g^\dagger g=I\)，型的相似变换则允许右端为非零标量；这个标量的实际取值范围依赖底域和型。Hermitian 型的共轭作用不能从公式中省略，复数结构则把酉群与实正交群联系起来。Clifford 代数将二次取值转为乘法关系；它的泛性质来自张量代数及双边理想商，具体维数结论仍须具体证明。
\end{chaptersummary}

\BookExercises

\begin{exercise}\label{b2:ex:10-left-right-orthogonals}
设 \(k\) 为域，在 \(k^2\) 上取 \(b(x,y)=x_1y_1+x_1y_2+x_2y_2\)。记 \(e_1,e_2\) 为标准基。证明型非退化，并对 \(W=ke_1\) 分别求 \(W^{\perp_r}\) 与 \({}^{\perp_l}W\)。说明为什么一般型不能只写一个没有侧别的正交补。
\end{exercise}
\begin{solution}
矩阵为 \(\begin{psmallmatrix}1&1\\0&1\end{psmallmatrix}\)，行列式为一，故非退化。\(b(e_1,y)=y_1+y_2\)，所以 \(W^{\perp_r}=k(1,-1)\)；而 \(b(x,e_1)=x_1\)，所以 \({}^{\perp_l}W=ke_2\)。两条直线不同，在特征二中也不同。非退化保证两边的维数正确，却不把不同变量中的正交条件变成同一个条件。
\end{solution}

\begin{exercise}\label{b2:ex:10-degenerate-orthogonals}
设 \(k\) 为域，\(V\) 为有限维 \(k\) 向量空间，\(b\) 为 \(V\) 上的对称型，允许退化，令 \(R=\operatorname{rad}b\)。证明对任意 \(W\le V\)，
\[
\dim W^\perp=\dim V-\dim W+\dim(W\cap R),
\qquad (W^\perp)^\perp=W+R.
\]
\end{exercise}
\begin{solution}
型在商空间 \(\overline V=V/R\) 上良定义：改变任一代表元所添的项都来自根空间，配对不变。所得型非退化，因为与全部商类正交的代表元正属于 \(R\)。记 \(\overline W=(W+R)/R\)，则 \(W^\perp/R=\overline W^\perp\)。由\cref{b2:thm:orthogonal-complement}，后者的维数为 \(\dim\overline V-\dim\overline W\)。代入 \(\dim\overline W=\dim W-\dim(W\cap R)\)，得到第一式。再次取正交补，商空间中得到 \(\overline W\)，取原像便是 \(W+R\)，证明第二式。
\end{solution}

\begin{exercise}\label{b2:ex:10-symplectic-basis}
在任意域 \(k\) 上，四维交替型的矩阵为
\[
B=\begin{pmatrix}0&1&1&0\\-1&0&0&0\\-1&0&0&1\\0&0&-1&0\end{pmatrix}.
\]
找出两个互相正交的辛平面，并给出辛基。说明该计算在特征二中是否仍有效。
\end{exercise}
\begin{solution}
第一对取 \(e_1,e_2\)。向量 \(e_3-e_2\) 与这两者的配对均为零，\(e_4\) 也与这两者正交，而 \(b(e_3-e_2,e_4)=1\)。所以 \(e_1,e_2,e_3-e_2,e_4\) 是按两维块排列的辛基；这四个向量线性无关，由原标准基到它们的变换可逆。只用了加减和单位系数，没有除以二，故在特征二中同样成立。
\end{solution}

\begin{exercise}\label{b2:ex:10-one-dimensional-square-class}
分类 \(\mathbb F_5\) 上全部非退化一维二次型，并从 \(\langle1\rangle,\langle2\rangle,\langle3\rangle,\langle4\rangle\) 中选出每个等距类的代表。再比较 \(\langle1\rangle\) 与 \(\langle2\rangle\) 在 \(\mathbb Q\) 和 \(\mathbb Q(\sqrt2)\) 上的等距性，说明扩大底域怎样影响平方类判据。
\end{exercise}
\begin{solution}
由\cref{b2:thm:quadratic-diagonalization}中的一维判据，只须分类非零系数的平方类。\(\mathbb F_5^\times\) 中的平方为 \(1,4\)，另一个平方类为 \(2,3\)，所以可取 \(\langle1\rangle\) 与 \(\langle2\rangle\) 作代表；\(\langle4\rangle\cong\langle1\rangle\)、\(\langle3\rangle\cong\langle2\rangle\)。

在有理数域上 \(2\) 不是平方：若互素整数满足 \(m^2=2n^2\)，则先得 \(m\) 偶，再得 \(n\) 偶，矛盾。因此两型不等距。扩大到 \(\mathbb Q(\sqrt2)\) 后，\(2=(\sqrt2)^2\)，两型等距。原域中不同的平方类可以在扩域后合并，原有等距关系则仍被保留。
\end{solution}

\begin{exercise}\label{b2:ex:10-diagonalization-radical}
在 \(\mathbb R^3\) 上，对
\[
q(x,y,z)=x^2+2xy+2y^2+2yz+z^2
\]
给出可逆坐标代换，使它对角化；求惯性、根空间，以及使限制型非退化的一个补空间。
\end{exercise}
\begin{solution}
令 \(u=x+y,v=y+z,w=y\)，逆变换为 \(x=u-w,y=w,z=v-w\)。于是 \(q=u^2+v^2\)，惯性为 \((2,0,1)\)。根空间为 \(u=v=0\)，即 \(\mathbb R(-1,1,-1)\)。取 \(w=0\)，得到 \(y=0\) 的平面，其限制型为 \(x^2+z^2\)，非退化；它与上述根空间构成直和。
\end{solution}

\begin{exercise}\label{b2:ex:10-inertia-parameter}
对实参数 \(t\)，分类 \(q_t(x,y,z)=x^2+2xy+ty^2-z^2\) 的惯性，并求出它何时非退化、何时具有非零各向同性向量。
\end{exercise}
\begin{solution}
写为 \((x+y)^2+(t-1)y^2-z^2\)。当 \(t>1,t=1,t<1\) 时，惯性依次为 \((2,1,0),(1,1,1),(1,2,0)\)。所以非退化恰在 \(t\ne1\) 时成立。无论 \(t\) 如何，\((1,0,1)\) 都是非零零点，故总有非零各向同性向量。
\end{solution}

\begin{exercise}\label{b2:ex:10-char-two-same-polar}
在 \(\mathbb F_2^2\) 上取 \(q_0(x,y)=xy\)、\(q_1(x,y)=x^2+xy+y^2\)。分别求保持 \(q_0\)、\(q_1\) 的线性自同构群，并与保持它们共同极化型的群比较。
\end{exercise}
\begin{solution}
\cref{b2:prop:quadratic-char-two-polar-failure}已经给出三条非零向量上的取值。保持 \(q_0\) 的变换必须保持零点集合 \(\{e_1,e_2\}\)，因而只能固定或交换这两个基向量；两者都保持 \(q_0\)，所得群为 \(\{I_2,\begin{psmallmatrix}0&1\\1&0\end{psmallmatrix}\}\cong C_2\)。\(q_1\) 在每个非零向量上都取一，故它由全部 \(\operatorname{GL}_2(\mathbb F_2)\) 保持。选择第一列有三种，第二列须与第一列线性无关，有两种，因此该群有六个元素；它忠实地置换三条非零向量，所以同构于 \(S_3\)。

共同的极化型矩阵为 \(J=\begin{psmallmatrix}0&1\\1&0\end{psmallmatrix}\)。二维计算给出 \(G^{\mathsf T}JG=(\det G)J\)，而 \(\mathbb F_2\) 中每个可逆矩阵的行列式都是一，故其保持群也是 \(\operatorname{GL}_2(\mathbb F_2)\)。保持极化型的变换未必保持原二次型，\(q_0\) 的群在这里是严格较小的子群。
\end{solution}

\begin{exercise}\label{b2:ex:10-binary-isotropy}
设 \(k\) 为特征不为二的域，\(d\in k^\times\)。证明 \(\langle1,-d\rangle\) 各向异性，当且仅当 \(d\) 不是平方；若 \(d\) 是平方，写出一组双曲基。
\end{exercise}
\begin{solution}
非零零点满足 \(x^2=dy^2\)。若 \(y=0\)，则 \(x=0\)，故必须 \(y\ne0\)，得到 \(d=(x/y)^2\)。反过来，若 \(d=c^2\)、\(c\ne0\)，取
\[
e=(c,1),\qquad f=(1/(2c),-1/(2c^2)).
\]
直接得到 \(q(e)=q(f)=0\) 及 \(b(e,f)=1\)，故为双曲基。
\end{solution}

\begin{exercise}\label{b2:ex:10-hyperbolic-split-explicit}
设 \(k\) 为特征不为二的域，在 \(k^3\) 上取 \(q(x,y,z)=2xy+z^2\)。以各向同性向量 \(e=(1,-1/2,1)\) 开始，取 \(f_0=(0,0,1)\)，再将它调整为与 \(e\) 配对为一的各向同性向量 \(f\)。实际完成一次双曲平面的分离，写出正交补的基和限制型。
\end{exercise}
\begin{solution}
相应对称型为 \(b\bigl((x,y,z),(u,v,w)\bigr)=xv+yu+zw\)。先有 \(b(e,f_0)=1\)、\(q(f_0)=1\)，于是按\cref{b2:lem:hyperbolic-plane-split}的修正方法取
\[
f=f_0-\frac{q(f_0)}2e=(-1/2,1/4,1/2).
\]
它满足 \(q(f)=0\)、\(b(e,f)=1\)。向量 \((x,y,z)\) 同时与 \(e,f\) 正交的条件为 \(-x+2y+2z=0\)、\(x-2y+2z=0\)。相加先得 \(z=0\)，再得 \(y=x/2\)，故正交补由 \(g=(1,1/2,0)\) 生成，且 \(q(g)=1\)。因此 \(e,f,g\) 给出 \(H\perp\langle1\rangle\) 的具体分解；这里调整 \(f_0\) 的系数由各向同性条件决定。
\end{solution}

\begin{exercise}\label{b2:ex:10-two-reflections}
在实三维型 \(q(x,y,z)=x^2+y^2-z^2\) 中，令 \(u=(1,0,0)\)、\(v=(1,1,1)\)。解释为何沿 \(u-v\) 的反射没有定义，并用两个有定义的反射把 \(u\) 送到 \(v\)。
\end{exercise}
\begin{solution}
\(q(u)=q(v)=1\)，但 \(q(u-v)=q(0,-1,-1)=0\)，反射公式的分母为零。另一方面 \(u+v=(2,1,1)\) 的二次取值为四，且 \(b(u,u+v)=2\)，故 \(\tau_{u+v}(u)=u-(u+v)=-v\)。\(q(v)=1\) 保证 \(\tau_v\) 有定义，并有 \(\tau_v(-v)=v\)。所求映射为 \(\tau_v\tau_{u+v}\)。
\end{solution}

\begin{exercise}\label{b2:ex:10-real-isotropic-dimension}
设有限维实向量空间上的二次型的惯性为 \((p_+,p_-,r)\)。证明其全各向同性子空间的最大维数为 \(\min(p_+,p_-)+r\)，并说明为什么根空间的贡献必须单列。
\end{exercise}
\begin{solution}
先在非退化标准型中讨论。若 \(T\) 全各向同性，它到正坐标空间和负坐标空间的两个投影都单射：其中一个投影为零，则一个非零向量的型值只能严格为正或严格为负，与零值矛盾。因此 \(\dim T\le\min(p_+,p_-)\)。把一条正坐标轴与一条负坐标轴配对，取各对的和向量，就得到达到上界的全各向同性子空间。

一般情形下，先除去根空间。\(T\) 的像全各向同性，故维数至多为 \(\min(p_+,p_-)\)，而核 \(T\cap\operatorname{rad}b\) 维数至多为 \(r\)。把前面构造的子空间与整个根空间直和即可达到总上界。根空间是零型，既不贡献正方向也不贡献负方向，所以它不是双曲分量，应另外加上其维数。
\end{solution}

\begin{exercise}\label{b2:ex:10-witt-binary-relation}
设 \(k\) 为特征不为二的域，\(a,b\in k\)，且 \(a,b,a+b\) 都非零。通过显式正交基证明
\[
\langle a,b\rangle\cong\left\langle a+b,\frac{ab}{a+b}\right\rangle.
\]
据此写出 \(W(k)\) 中的一条关系。
\end{exercise}
\begin{solution}
在原正交基 \(e,f\) 中取 \(u=e+f\)、\(v=(be-af)/(a+b)\)。两者的配对为 \((ab-ab)/(a+b)=0\)，且 \(q(u)=a+b\)、\(q(v)=(ab^2+a^2b)/(a+b)^2=ab/(a+b)\)。换基矩阵的行列式为负一，所以是基。这给出等距，并在 Witt 群中得到
\[
\bigl[\langle a\rangle\bigr]+\bigl[\langle b\rangle\bigr]
=\bigl[\langle a+b\rangle\bigr]+\left[\left\langle\frac{ab}{a+b}\right\rangle\right].
\]
最后两层定界符只用于区分一维型与它的 Witt 类。
\end{solution}

\begin{exercise}\label{b2:ex:10-real-complex-witt-map}
证明将实二次型的系数看成复数，诱导群同态 \(W(\mathbb R)\to W(\mathbb C)\)。在\cref{b2:prop:witt-real-complex}的识别下求出该同态及其核。
\end{exercise}
\begin{solution}
实等距矩阵在复数域仍可逆并满足同一个合同等式，双曲平面扩域后仍是双曲平面，所以 Witt 等价及类的加法都被保持。实惯性为 \((p_+,p_-,0)\) 的型，复维数为 \(p_++p_-\)，其奇偶性等于 \(p_+-p_-\) 的奇偶性。因此该同态就是整数模二，核为 \(2\mathbb Z\)。
\end{solution}

\begin{exercise}\label{b2:ex:10-orthogonal-one-dimension}
设 \(k\) 为特征不为二的域。求 \(k\) 上一维非退化二次型的正交群和特殊正交群。再设二维双曲平面的矩阵为 \(B_H=\begin{psmallmatrix}0&1\\1&0\end{psmallmatrix}\)，证明 \(\operatorname{diag}(t,t^{-1})\)、\(t\in k^\times\)，都在其特殊正交群中。
\end{exercise}
\begin{solution}
一维线性自同构为乘 \(c\ne0\)，保持型要求 \(c^2=1\)，故正交群为 \(\{1,-1\}\)，特殊正交群为 \(\{1\}\)。双曲平面上，变换 \((x,y)\mapsto(tx,t^{-1}y)\) 保持 \(2xy\)，且行列式为一，故满足所述条件。这个一参数族表明不定型的保持群允许沿两条各向同性轴作相反的伸缩。
\end{solution}

\begin{exercise}\label{b2:ex:10-first-order-isometries}
设 \(k\) 为域，\(n\ge1\)，\(R=k[\varepsilon]/(\varepsilon^2)\)，\(B\in M_n(k)\) 可逆，给定 \(X\in M_n(k)\)。证明 \(G=I_n+\varepsilon X\) 在 \(R\) 上保持 \(B\)，当且仅当 \(X^{\mathsf T}B+BX=0\)。\(\operatorname{char}k\ne2\) 时，分别计算 \(B=I_n\) 与 \(B=J_m\)（此时 \(n=2m\)）的解空间维数。
\end{exercise}
\begin{solution}
双数中的 \(\varepsilon^2=0\) 使保持型的等式只留下相对于 \(\varepsilon\) 的一次变化，因此所得条件将是关于 \(X\) 的线性方程。\(G\) 的逆为 \(I_n-\varepsilon X\)。展开 \(G^{\mathsf T}BG\)，得到 \(B+\varepsilon(X^{\mathsf T}B+BX)\)。由于 \(1,\varepsilon\) 在 \(k\) 上线性无关，保持 \(B\) 恰好等价于所述等式。\(B=I_n\) 时，\(X\) 反对称且对角元为零，独立参数为上三角部分的 \(n(n-1)/2\) 项。

若 \(B=J_m\)，写 \(X=\begin{psmallmatrix}A&C\\D&E\end{psmallmatrix}\)。分块相乘给出 \(E=-A^{\mathsf T}\)、\(C=C^{\mathsf T}\)、\(D=D^{\mathsf T}\)。所以维数为 \(m^2+2\cdot m(m+1)/2=m(2m+1)\)。
\end{solution}

\begin{exercise}\label{b2:ex:10-indefinite-hermitian}
在 \(\mathbb C^2\) 上取 Hermitian 矩阵 \(H=\begin{psmallmatrix}1&1\\1&0\end{psmallmatrix}\)。求其惯性并给出 \(P^*HP\) 的标准形。再证明
\[
A=\begin{pmatrix}5/3&4/3\\4/3&5/3\end{pmatrix}
\]
属于 \(\operatorname U(1,1)\)，却不属于 \(\operatorname U(2)\)。
\end{exercise}
\begin{solution}
对角取值为 \(|z_1+z_2|^2-|z_2|^2\)。取 \(P=\begin{psmallmatrix}1&-1\\0&1\end{psmallmatrix}\)，即 \(z_1=u-v,z_2=v\)，得到 \(P^*HP=\operatorname{diag}(1,-1)\)，惯性为 \((1,1,0)\)。

记 \(c=5/3,s=4/3\)，有 \(c^2-s^2=1\)。相乘得 \(A^{\mathsf T}\operatorname{diag}(1,-1)A=\operatorname{diag}(1,-1)\)，故 \(A\in\operatorname U(1,1)\)。但第一列的普通模平方为 \(c^2+s^2=41/9\ne1\)，所以 \(A^*A\ne I\)，不在标准酉群中。
\end{solution}

\begin{exercise}\label{b2:ex:10-unitary-real-blocks}
设 \(n\ge1\)，在 \(\mathbb R^{2n}\) 上令 \(J=\begin{psmallmatrix}0&-I_n\\I_n&0\end{psmallmatrix}\)。求所有与 \(J\) 交换的实矩阵的分块形式，并证明其中的正交矩阵恰对应复酉矩阵。
\end{exercise}
\begin{solution}
将四个 \(n\) 阶块代入 \(TJ=JT\)，得到且仅得到
\[
T=\begin{pmatrix}A&-B\\B&A\end{pmatrix}.
\]
它表示复矩阵 \(A+iB\) 的实线性作用。\(T^{\mathsf T}T=I\) 等价于 \(A^{\mathsf T}A+B^{\mathsf T}B=I\) 与 \(A^{\mathsf T}B=B^{\mathsf T}A\)，而这两式恰是 \((A+iB)^*(A+iB)=I\) 的实部和虚部。因此得到所述对应。
\end{solution}

\begin{exercise}\label{b2:ex:10-adjoint-image-kernel}
设 \(k\) 为特征不为二的域，在 \(k^2\) 上取矩阵为 \(H=\begin{psmallmatrix}0&1\\1&0\end{psmallmatrix}\) 的对称双线性型。证明 \(T=\begin{psmallmatrix}0&1\\0&0\end{psmallmatrix}\) 关于此型自伴随，且 \(T^2=0\) 而 \(T\ne0\)；求 \(\operatorname{im}T\)、\(\ker T\) 及 \((\operatorname{im}T)^\perp\)。再证明有限维实正定内积空间或复正定 Hermitian 空间上的自伴随算子若满足 \(T^2=0\)，就必为零，并指出两种情形的区别。
\end{exercise}
\begin{solution}
\(H^{-1}=H\)，相乘得到 \(H^{-1}T^{\mathsf T}H=T\)，故由\cref{b2:prop:form-adjoint-operator}，\(T^\dagger=T\)。又 \(T(e_2)=e_1\)、\(T(e_1)=0\)，所以 \(T^2=0\ne T\)，且 \(\operatorname{im}T=\ker T=ke_1\)。条件 \(b(e_1,v)=v_2=0\) 给出 \((\operatorname{im}T)^\perp=ke_1\)，也与正文的像、核恒等式一致。像与核在这里相等，它们并不构成直和。

在正定情形，自伴随性与 \(T^2=0\) 给出
\[
b(Tv,Tv)=b(v,T^\dagger Tv)=b(v,T^2v)=0.
\]
正定性迫使每个 \(Tv\) 为零，故 \(T=0\)。前一例的 \(ke_1\) 是全各向同性直线，非零向量也可有零取值；非退化性本身不能替代正定性。
\end{solution}

\begin{exercise}\label{b2:ex:10-indefinite-similitude-multipliers}
设 \(b\) 是 \(\mathbb R^2\) 上矩阵为 \(B=\operatorname{diag}(1,-1)\) 的双线性型。取 \(P=\begin{psmallmatrix}1&1\\0&1\end{psmallmatrix}\)，计算换基后的矩阵 \(P^{\mathsf T}BP\)，并判断 \(P\) 是否保持 \(b\)，或是否将它变为非零标量倍数。再求 \(b\) 的相似乘子的像，给出一个乘子为负一的变换，并说明正定型为何不可能有负乘子。
\end{exercise}
\begin{solution}
直接相乘得 \(P^{\mathsf T}BP=\begin{psmallmatrix}1&1\\1&0\end{psmallmatrix}\)。它既不等于 \(B\)，也不是 \(B\) 的标量倍数，所以 \(P\) 虽可用于换基，却不是 \(b\) 的等距变换或型的相似变换。交换两个坐标的矩阵 \(S=\begin{psmallmatrix}0&1\\1&0\end{psmallmatrix}\) 满足 \(S^{\mathsf T}BS=-B\)，是乘子为负一的型的相似变换，却不是等距变换。每个正数 \(c\) 由 \(\sqrt c I_2\) 实现，每个负数 \(c\) 由 \(\sqrt{-c}S\) 实现，故乘子的像为 \(\mathbb R^\times\)。若某个变换 \(g\) 对正定型 \(h\) 有负乘子 \(c\)，则非零 \(v\) 满足 \(h(gv,gv)=c\cdot h(v,v)<0\)，与正定性矛盾。
\end{solution}

\begin{optionalexercise}\label{b2:ex:10-clifford-line}
求实一维二次型的三种 Clifford 代数模型 \(\mathbb R\times\mathbb R\)、\(\mathbb C\)、\(D=\mathbb R[\varepsilon]/(\varepsilon^2)\) 的全部含幺 \(\mathbb R\) 代数自同构。分别求出三个代数中的幂等元，说明它们作为实向量空间维数相同却不能同构。
\end{optionalexercise}
\begin{solution}
\cref{b2:prop:clifford-line-models}给出这三个模型。乘积代数的幂等元为 \((0,0),(1,0),(0,1),(1,1)\)。自同构固定零与单位，故必须固定或交换两个非平凡幂等元；它们又线性生成该代数，所以自同构只有恒等与交换两因子。\(\mathbb C\) 中 \(i\) 的像须满足平方为负一，只能为 \(i\) 或 \(-i\)，两者分别给出恒等与复共轭。

在 \(D\) 中，\(1,\varepsilon\) 为实基。若 \((a+b\varepsilon)^2=0\)，常数项给出 \(a=0\)，故 \(\varepsilon\) 在任何自同构下必变为 \(b\varepsilon\)。可逆性等价于 \(b\ne0\)，而每个这样的缩放确实给出自同构，因此 \(\operatorname{Aut}_{\mathbb R}(D)\cong\mathbb R^\times\)。

\(\mathbb C\) 中幂等元只有零与一。\(D\) 中的幂等条件为 \(a^2=a\)、\((2a-1)b=0\)，同样只有零与一。非平凡幂等元将乘积代数与后两者区分；后两者又由非零幂零元的存在性区分：\(\varepsilon\ne0\) 但 \(\varepsilon^2=0\)，而域 \(\mathbb C\) 没有非零幂零元。所以三者虽都是二维实代数，却两两不同构。
\end{solution}

\begin{optionalexercise}\label{b2:ex:10-clifford-char-two}
在 \(\mathbb F_2^2\) 上取 \(q(x,y)=x^2+y^2\)。证明其 Clifford 代数同构于
\[
\mathbb F_2[u,v]/(u^2,v^2).
\]
直接给出该代数的一个向量空间基，不引用一般 Clifford 代数的维数结论。
\end{optionalexercise}
\begin{solution}
令两个标准基的像为 \(e,f\)，则 \(e^2=f^2=1\)。极化型为零，由\eqref{b2:eq:clifford-polar-relation}，\(ef+fe=0\)，在特征二中即 \(ef=fe\)。置 \(u=e+1,v=f+1\)，得到 \(u^2=v^2=0\)，且它们交换，因此从题中商代数有一个满同态到 Clifford 代数。

反过来，将 \((x,y)\) 送到 \(x(1+u)+y(1+v)\)。因交换性与特征二，这个元素的平方为 \(x^2+y^2\)，由泛性质得到反向同态。两复合固定各生成元，故为同构。多项式按次数分别对 \(u^2,v^2\) 取余，唯一得到 \(a+bu+cv+duv\)：理想 \((u^2,v^2)\) 中每个单项式至少有一个指数不小于二，不能消去这四个单项式。因此 \(1,u,v,uv\) 是基，维数为四。
\end{solution}
